% concatenated from article7.tex, preamble.tex, font.tex, titlepage.tex
\documentclass[11pt,twoside]{article}

\newcommand{\HeadTitle}{Compact Coefficient Formulae for Logarithmic Tangent and Hyperbolic Integrals}

\newcommand{\HeadTitleTwo}{\begin{center}
\Large{\textit{Compact Coefficient Formulae for Logarithmic Tangent and Hyperbolic Integrals}}
\end{center}}

% =========================
% PREAMBLE COMMUN (FIGÉ)
% =========================

% --- Géométrie & maths ---
\usepackage[margin=0.8in]{geometry}
\usepackage{amsmath, amssymb, amsfonts, mathtools}
\usepackage{amsthm}
\usepackage{mathrsfs}
\usepackage{stmaryrd}
\usepackage{esint}
\usepackage{eqnalign}

% --- Graphiques & tableaux ---
\usepackage{graphicx}
\usepackage{enumitem}
\usepackage[most]{tcolorbox}
\usepackage{float}
\usepackage{listings}
\usepackage{array}
\usepackage{booktabs}
\usepackage{xcolor}
\usepackage{emptypage}
\usepackage{tensor}
\usepackage{tabularx}

% --- En-têtes & références ---
\usepackage{fancyhdr}
\usepackage{hyperref}
\usepackage[nameinlink,noabbrev]{cleveref}

% --- Titres & fontes ---
\usepackage{titlesec}
\usepackage[T1]{fontenc}
\usepackage[utf8]{inputenc}
\usepackage{lmodern}
\usepackage{titling}

% --- Bibliographie ---
\usepackage[backend=biber,style=numeric]{biblatex}
\addbibresource{literatur.bib}
\setlength\bibitemsep{1em}

% --- Commandes perso ---
\newcommand{\bigint}{
  \mathop{\scalebox{2.5}{$\displaystyle\int$}}\limits
}
\newcommand{\EulerB}[2]{\left\langle {#1 \atop #2} \right\rangle^{\!B}}
\newcommand{\Arg}{\operatorname{Arg}}
\DeclareMathOperator{\sech}{sech}
\DeclareMathOperator*{\Res}{Res}
% =========================
% MÉTADONNÉES (IMPORTANT)
% =========================
\providecommand{\HeadTitle}{} % <-- TITRE VARIABLE
\providecommand{\HeadTitleTwo}{} % <-- TITRE VARIABLE
\providecommand{\HeadAuthor}{Luc Ramsès TALLA WAFFO} % <-- FIXE

% =========================
% FANCYHDR
% =========================
\pagestyle{fancy}
\fancyhf{}
\fancyhead[CO]{\small\itshape\HeadTitle}
\fancyhead[CE]{\small\itshape\HeadAuthor}
\fancyfoot[C]{\thepage}

% =========================
% TITRES DE SECTIONS
% =========================
\titleformat{\section}[block]
  {\normalfont\large\bfseries\itshape\centering}
  {§\thesection.}
  {1em}
  {}

\titleformat{\subsection}
  {\normalfont\normalsize\bfseries}
  {\thesubsection}
  {1em}
  {}
	
	\newcommand{\K}{\mathbb{K}}
\newcommand{\R}{\mathbb{R}}
\newcommand{\C}{\mathbb{C}}
\newcommand{\lc}{\operatorname{lc}}
\newcommand{\End}{\operatorname{End}}
\newcommand{\Span}{\operatorname{span}}
\newcommand{\degzero}{-\infty}

% =========================
% ENVIRONNEMENTS THÉORÈMES
% =========================
\newtheorem{theorem}{Theorem}[section]
\newtheorem{lemma}[theorem]{Lemma}
\newtheorem{proposition}[theorem]{Proposition}
\newtheorem{corollary}[theorem]{Corollary}
\newtheorem{conjecture}[theorem]{Conjecture}
\newtheorem{remark}[theorem]{Remark}
\newtheorem{definition}[theorem]{Definition}
\newtheorem{example}[theorem]{Example}

% =========================
% CLEVERREF
% =========================
\crefname{example}{example}{examples}
\Crefname{example}{Example}{Examples}

\crefname{corollary}{corollary}{corollaries}
\Crefname{corollary}{Corollary}{Corollaries}

\crefname{definition}{definition}{definitions}
\Crefname{definition}{Definition}{Definitions}

\crefname{remark}{remark}{remarks}
\Crefname{remark}{Remark}{Remarks}

\crefname{conjecture}{conjecture}{conjectures}
\Crefname{conjecture}{Conjecture}{Conjectures}

\crefname{lemma}{lemma}{lemmas}
\Crefname{lemma}{Lemma}{Lemmas}

\crefname{proposition}{proposition}{propositions}
\Crefname{proposition}{Proposition}{Propositions}

\crefname{theorem}{theorem}{theorems}
\Crefname{theorem}{Theorem}{Theorems}

\numberwithin{equation}{section}
% Polices système et Unicode
%\usepackage{fontspec}
%\setmainfont{Times New Roman}    % Texte principal
%\setsansfont{Arial}              % Optionnel : sans-serif
%\setmonofont{Courier New}        % Optionnel : monospace

% Maths compatibles Times
%\usepackage{unicode-math}
%\usepackage{lmodern}
%\setmathfont{TeX Gyre Termes Math}  % Ou TeX Gyre Termes Math si tu veux
%\setmathfont{Latin Modern Math}  % Ou TeX Gyre Termes Math si tu veux

\usepackage{fontspec}
\setmainfont{texgyretermes-regular.otf}[
  BoldFont = texgyretermes-bold.otf,
  ItalicFont = texgyretermes-italic.otf,
  BoldItalicFont = texgyretermes-bolditalic.otf
]

\begin{document}

\thispagestyle{fancy}

\vspace{0.2cm}

\begin{center}
\Large{\HeadTitleTwo}
\end{center}

\hspace{3cm}

\begin{center}
Luc Ramsès TALLA WAFFO \\
Technische Universität Darmstadt\\
Karolinenplatz 5, 64289 Darmstadt, Germany\\
ramses.talla@stud.tu-darmstadt.de\\
\vspace{0.5cm}
\today
\end{center}


\begin{abstract}
We develop compact coefficient-extraction formulae for several families of
hyperbolic, logarithmic tangent, and Malmsten-type integrals whose values are
finite linear combinations of odd zeta values and even Dirichlet beta values.
The principal advantage of these formulae is that coefficients previously
encoded by recursive arrays or nested finite sums are replaced by a single
coefficient of an explicit elementary expression.  This makes the
coefficients easier to compute, keeps the dependence on the parameters
visible, and reveals structural features---such as vanishing ranges,
extremal coefficients, and sign patterns---without hidden cancellations.
For shifted hyperbolic integrals with numerator $\sinh((2k+1)x)$, the
coefficients are expressed through Chebyshev--arcsine extractions.  The same
mechanism yields Laurent coefficient formulae for logarithmic tangent
integrals and leads to direct proofs of simple initial and terminal
coefficients, including a parity-free terminal identity.  For $m,n\geq1$,
$m\geq n$, and $m+n$ even, we prove
\[
\int_0^\infty\frac{\tanh^{m+1}x}{x^{n+1}}\,dx
=
(-1)^{(m-n)/2}
\sum_{p=\lceil n/2\rceil}^{(m+n)/2}
\binom{2p}{n}
(2^{2p+1}-1)
\frac{\zeta(2p+1)}{\pi^{2p}}
[u^{m+n-2p}](u\cot u)^{m+1}.
\]
In the diagonal case, this gives the family
$\int_0^\infty(\tanh x/x)^N\,dx$ in a direct, non-recursive form and
makes the disappearance of the initial zeta values immediate.
\end{abstract}

\vspace{0.2cm}

\paragraph{Notation.}
We write \(\mathbb N=\{1,2,\ldots\}\).  The Riemann zeta function and the
Dirichlet beta function are denoted by \(\zeta(s)\) and \(\beta(s)\),
respectively.  The Chebyshev polynomial of the second kind \cite{MasonHandscomb2003}
\cite{Mason1993}\cite{Rivlin1990}
\cite{NIST:DLMF} is denoted by
\(U_m\) and is normalized by
\[
U_m(\cos\theta)=\frac{\sin((m+1)\theta)}{\sin\theta}.
\]
For a power series or Laurent series \(F\), the notation \([x^r]F(x)\)
denotes the coefficient of \(x^r\) in the expansion at the origin.

\vspace{0.5cm}

\section*{Introduction}

Definite integrals involving hyperbolic functions and logarithmic tangent
kernels often admit finite evaluations in terms of odd values of the Riemann
zeta function and even values of the Dirichlet beta function.  In many such
formulae, however, the special values themselves are not the main obstacle:
the essential complexity is carried by the accompanying coefficients.  These
coefficients are frequently presented through recurrences, nested finite sums,
or auxiliary arrays, which may be effective for computation at fixed
parameters but tend to conceal the dependence on those parameters and the
structure of the resulting expansion.  Kyrion \cite{Kyrion2025}, for example,
studied integrals of the form
\[
\int_0^\infty \frac{\tanh x}{x}\sech^L x\,e^{-Tx}\,dx
\]
and, in the case \(T=0\), obtained recursive coefficient formulae for the
zeta and beta expansions of
\[
\int_0^\infty \frac{\tanh x}{x}\sech^{2N}x\,dx
\qquad\text{and}\qquad
\int_0^\infty \frac{\tanh x}{x}\sech^{2N+1}x\,dx.
\]
Related contour evaluations for shifted hyperbolic integrals and logarithmic
tangent integrals were obtained in \cite{talla_waffo_integral_2025}.

\vspace{0.1cm}

The aim of the present paper is to recast coefficients of this type as
explicit coefficient extractions from elementary functions.  The guiding
principle is simple: instead of generating a family of coefficients
recursively, one identifies the entire family with coefficients of a single
power series or Laurent series.  In the cases considered here, the relevant
expressions involve Chebyshev polynomials of the second kind, powers of
\(\arcsin x\), and the generating functions \((u\cot u)^r\) and
\((w/\sinh w)^r\).  This reformulation has several advantages.  It gives
direct finite formulae, keeps the parameter dependence visible, and makes
structural properties such as support, vanishing ranges, endpoint
coefficients, and sign patterns accessible without resolving hidden
cancellations in nested sums.  The Chebyshev identities used throughout are
standard; see, for example, \cite[Section~4.6]{MasonHandscomb2003},
\cite{Mason1993,Rivlin1990}, and \cite[Chapter~18]{NIST:DLMF}.

\vspace{0.1cm}

We begin with shifted hyperbolic integrals having odd numerator frequency,
\[
\int_{0}^{\infty}
\frac{\sinh((2k+1)x)}{x\cosh^{2n+1}x}\,dx
\qquad\text{and}\qquad
\int_{0}^{\infty}
\frac{\sinh((2k+1)x)}{x\cosh^{2n}x}\,dx.
\]
The zeta and beta coefficients in these families are reduced to the single
extractions
\[
[x^{2n}]\,U_{2k}(x)(\arcsin x)^{2p}
\qquad\text{and}\qquad
[x^{2n-1}]\,U_{2k}(x)(\arcsin x)^{2p-1},
\]
respectively.  The case \(k=0\) recovers the unshifted families considered by
Kyrion, whereas arbitrary \(k\) is incorporated by the same Chebyshev factor.
Linearization of odd powers of \(\sinh x\) then collapses the corresponding
Chebyshev combination to \((x^2-1)^k\), yielding compact formulae for
integrals with numerator \(\sinh^{2k+1}x\).  As a consequence, two families of
integrals involving \(1/\operatorname{arctanh}x\) on \((0,1)\) follow by the
substitution \(x=\tanh t\).

\vspace{0.1cm}

A parallel argument treats the even-frequency numerator \(\sinh(2kx)\).  By
combining its binomial expansion with a Chebyshev polynomial identity, the
corresponding integrals with even and odd powers of \(\cosh x\) in the
denominator are again expressed by one coefficient extraction, now involving
\(U_{2k-1}\).  Thus the odd- and even-frequency shifted families fit into the
same coefficient-extraction framework, although the parity of the denominator
determines whether zeta or beta values occur.

\vspace{0.1cm}

The same viewpoint applies to the logarithmic tangent integrals
\[
\int_0^{\pi/4}\frac{\sin(4nx)}{\ln(\tan x)}\,dx
\qquad\text{and}\qquad
\int_0^{\pi/4}\frac{\cos((4n-2)x)}{\ln(\tan x)}\,dx.
\]
After combining multiple-angle formulae with binomial identities for
Chebyshev polynomials, every zeta or beta coefficient becomes the coefficient
of \(x\) in a single Laurent product involving
\(U_{2n-1}(1/x)\) or \(U_{2n-2}(1/x)\) and a power of \(\arcsin x\).  In
addition to providing the complete expansions in compact form, these Laurent
representations isolate the boundary coefficients.  In particular, they give
a direct coefficient-extraction proof for the first beta coefficient and lead
to the parity-independent terminal identity
\[
[x]U_{r-1}\!\left(\frac1x\right)(\arcsin x)^r=2^{r-1}.
\]

\vspace{0.1cm}

The coefficient method also extends beyond the preceding contour families.  For
\(m,n\geq1\), \(m\geq n\), and \(m+n\) even, we consider
\[
\int_0^\infty\frac{\tanh^{m+1}x}{x^{n+1}}\,dx.
\]
Repeated integration by parts reduces this integral to a finite combination
of the odd-sinh family.  An auxiliary polynomial describing repeated
differentiation of \(\tanh^{m+1}x\), together with a cotangent derivative
identity, then collapses the remaining finite sum to a single coefficient of
\((u\cot u)^{m+1}\).  The resulting formula is
\[
\int_0^\infty
\frac{\tanh^{m+1}x}{x^{n+1}}\,dx
=
(-1)^{(m-n)/2}
\sum_{p=\lceil n/2\rceil}^{(m+n)/2}
\binom{2p}{n}
(2^{2p+1}-1)
\frac{\zeta(2p+1)}{\pi^{2p}}
[u^{m+n-2p}](u\cot u)^{m+1}.
\]
This formula is finite, explicit, and non-recursive.  Its lower summation bound
makes the disappearance of the initial zeta values immediate, while the
constant term of \((u\cot u)^{m+1}\) gives the highest zeta coefficient
directly.  The same representation also yields a simple first coefficient
when \(n\) is even and proves alternation of the nonzero zeta coefficients in
the relevant range.  In the diagonal case \(m=n\), it provides 
\[
\int_0^\infty
\left(\frac{\tanh x}{x}\right)^{n+1}\,dx
=
\sum_{p=\lceil n/2\rceil}^{n}
\binom{2p}{n}
(2^{2p+1}-1)
\frac{\zeta(2p+1)}{\pi^{2p}}
[u^{2n-2p}](u\cot u)^{n+1},
\]
replacing the nested and recursive coefficients that
occur in earlier descriptions of the same integral given by Kyrion
\cite{Kyrion2025}.
\vspace{0.1cm}
After suitable choice of parameters $m$ and $n$, we also deduce the closed-form of the integral$\displaystyle\int_0^\infty
\frac{\tanh^{2n}x}{x^2}\,dx$ which occurs in \cite{TallaWaffo2025arxiv2511.02843}
\[
\int_0^\infty
\frac{\tanh^{2n}x}{x^2}\,dx
=
(-1)^{n-1}
\sum_{p=1}^{n}
2p\,(2^{2p+1}-1)
\frac{\zeta(2p+1)}{\pi^{2p}}
[u^{2n-2p}](u\cot u)^{2n}
\]
Finally, the coefficient-extraction approach is applied to Blagouchine's
logarithmic hyperbolic integrals
\[
\mathcal I_r(a,b)
=
\int_0^\infty
\frac{\ln(x^2+a^2)}{\cosh^r(bx)}\,dx.
\]
Blagouchine \cite{Blagouchine2014} gave parity-dependent general forms and
tabulated the rational coefficients appearing in them.  Here those
coefficients are made explicit for all indices: both coefficient families are
obtained directly from
\[
\left(\frac{w}{\sinh w}\right)^r.
\]
This gives closed coefficient-extraction formulae for the even and odd cases
within the same elementary generating-function framework used throughout the
paper.

\vspace{0.1cm}

A final application concerns Malmsten-type integrals for the even beta values
$\beta(2n)$ and the odd zeta values $\zeta(2n+1)$.  Starting from finite
Eulerian expansions, we obtain a rigorous proof of the compact
polylogarithmic representations previously derived in
\cite{TallaWaffo2025arxiv2511.02843}, without relying on an interchange of an
infinite series with the integral.  The central mechanism is an
Eulerian--Chebyshev coefficient collapse.  If
$\left\langle {m\atop k}\right\rangle$ and
$\left\langle {m\atop k}\right\rangle^{B}$ denote Eulerian numbers of type
$A$ and type $B$, respectively, then the two finite Chebyshev sums satisfy
\[
\begin{dcases}
\displaystyle
[z^{2n-1}](\arcsin z)^{2p-1}
\sum_{k=0}^{n-1}(-1)^k
\left\langle {2n-1\atop k}\right\rangle^{B}
U_{2n-2k-2}(z)
=
(-1)^{n-1}2^{2n-2}(2n-1)!\,\delta_{p,n},
\\[3ex]
\displaystyle
[z^{2n}](\arcsin z)^{2p}
\sum_{k=0}^{n-1}(-1)^k
\left\langle {2n\atop k}\right\rangle
U_{2n-2k-2}(z)
=
\frac{(-1)^{n-1}}{2}(2n)!\,\delta_{p,n}.
\end{dcases}
\]
Thus every off-diagonal term $p\neq n$ vanishes.  For $p>n$ this is forced by
the order at the origin, whereas for $p<n$ it results from an exact
cancellation among nonzero contributions.  After normalization, these
identities may be viewed as a triangular biorthogonality relation between
powers of $\arcsin z$ and the corresponding Eulerian--Chebyshev polynomial
families.  This collapse is precisely what reduces the finite Eulerian
expansions to a single beta or zeta value.

\vspace{0.1cm}

The paper is organized as follows.  In \cref{sec:coefficient-formulae}, we
convert the coefficient sums for odd-frequency shifted hyperbolic integrals
into Chebyshev--arcsine extractions, derive the corresponding zeta and beta
formulae, pass to odd powers of \(\sinh x\), and obtain the reciprocal
hyperbolic arctangent corollaries.  In
\cref{sec:coefficient-formulae_double_angle}, we treat the even-frequency
family \(\sinh(2kx)\).  In
\cref{sec:coefficient-formulae-logarithmic-tangent-integrals}, we derive the
compact Laurent coefficient formulae for the logarithmic tangent integrals and
prove the first and terminal coefficient identities.  In
\cref{section:general-tangent-integral}, we develop the auxiliary polynomial,
perform the repeated integration-by-parts reduction, prove the cotangent
collapse, and derive the general tangent-power formula together with its
structural consequences and examples.  In
\cref{sec:blagouchine-coefficients}, we obtain explicit coefficient formulae
for Blagouchine's even- and odd-power logarithmic hyperbolic integral
families.  Finally, in \cref{sec:Malmsten-integrals}, we give the finite
Eulerian proof of the compact Malmsten-type polylogarithmic representations,
establish the Eulerian--Chebyshev coefficient collapse, and derive the
Kronecker-delta biorthogonality relation underlying the cancellation.

\vspace{0.3cm}

\section{Coefficient Formulae for $\displaystyle\int_{0}^{\infty}
\frac{\sinh((2k+1)x)}{x\cosh^{n}x}\,dx$ and $\displaystyle\int_{0}^{\infty}
\frac{\sinh^{2k+1}x}{x\cosh^{n}x}\,dx$}\label[section]{sec:coefficient-formulae}


\begin{lemma}[Conversion of the \(c\)- and \(d\)-sums]
\label[lemma]{lemma:cd-sums-to-arcsine-chebyshev}
Let \(n\in\mathbb N\), \(0\le k<n\), and \(1\le p\le n\). Let
\[
\left(\frac{t}{\sinh t}\right)^{2n+1}
=
\sum_{j\ge0}c_{j,n}t^j,
\qquad
\left(\frac{t}{\sinh t}\right)^{2n}
=
\sum_{j\ge0}d_{j,n}t^j.
\]
Let \(U_m\) denote the Chebyshev polynomial of the second kind. Then
\begin{equation}\label{eq:c-sum-to-arcsine-chebyshev}
\sum_{m=p}^{n}
c_{2n-2m,n}
\frac{(2k+1)^{2m-2p}}{(2m-2p)!}
=
(-1)^{n+p+k}
[x^{2n}]\,U_{2k}(x)(\arcsin x)^{2p},
\end{equation}
and
\begin{equation}\label{eq:d-sum-to-arcsine-chebyshev}
\sum_{m=0}^{n-p}
d_{2m,n}
\frac{(2k+1)^{2n-2m-2p}}
{(2n-2m-2p)!}
=
(-1)^{n+p+k}
[x^{2n-1}]\,U_{2k}(x)(\arcsin x)^{2p-1}.
\end{equation}
\end{lemma}

\begin{proof}
By the Cauchy product with
\[
\cosh((2k+1)t)
=
\sum_{\ell\ge0}
\frac{(2k+1)^{2\ell}t^{2\ell}}{(2\ell)!},
\]
we obtain
\begin{equation}\label{eq:cd-cauchy-product-local}
\begin{cases}
\displaystyle
\sum_{m=p}^{n}
c_{2n-2m,n}
\frac{(2k+1)^{2m-2p}}{(2m-2p)!}
=
[t^{2n-2p}]
\left(\frac{t}{\sinh t}\right)^{2n+1}
\cosh((2k+1)t),
\\[2.6em]
\displaystyle
\sum_{m=0}^{n-p}
d_{2m,n}
\frac{(2k+1)^{2n-2m-2p}}
{(2n-2m-2p)!}
=
[t^{2n-2p}]
\left(\frac{t}{\sinh t}\right)^{2n}
\cosh((2k+1)t).
\end{cases}
\end{equation}
By Cauchy's coefficient formula and the substitution \(t=iz\),
\eqref{eq:cd-cauchy-product-local} becomes
\begin{equation}\label{eq:cd-contour-local}
\begin{cases}
\displaystyle
\sum_{m=p}^{n}
c_{2n-2m,n}
\frac{(2k+1)^{2m-2p}}{(2m-2p)!}
=
(-1)^{n+p}
\frac{1}{2\pi i}
\oint_{\Gamma}
\left(\frac{z}{\sin z}\right)^{2n+1}
\cos((2k+1)z)
\frac{dz}{z^{2n-2p+1}},
\\[3.0em]
\displaystyle
\sum_{m=0}^{n-p}
d_{2m,n}
\frac{(2k+1)^{2n-2m-2p}}
{(2n-2m-2p)!}
=
(-1)^{n+p}
\frac{1}{2\pi i}
\oint_{\Gamma}
\left(\frac{z}{\sin z}\right)^{2n}
\cos((2k+1)z)
\frac{dz}{z^{2n-2p+1}},
\end{cases}
\end{equation}
where \(\Gamma\) is a sufficiently small positively oriented circle around
\(0\).

Now put
\[
x=\sin z.
\]
Then \(z=\arcsin x\) and \(dz=dx/\sqrt{1-x^2}\). Hence
\eqref{eq:cd-contour-local} becomes
\begin{equation}\label{eq:cd-after-sine-substitution-local}
\begin{cases}
\displaystyle
\sum_{m=p}^{n}
c_{2n-2m,n}
\frac{(2k+1)^{2m-2p}}{(2m-2p)!}
=
(-1)^{n+p}
\frac{1}{2\pi i}
\oint_{\sin\Gamma}
\frac{\cos((2k+1)\arcsin x)}{\sqrt{1-x^2}}
\frac{(\arcsin x)^{2p}}{x^{2n+1}}\,dx,
\\[3.0em]
\displaystyle
\sum_{m=0}^{n-p}
d_{2m,n}
\frac{(2k+1)^{2n-2m-2p}}
{(2n-2m-2p)!}
=
(-1)^{n+p}
\frac{1}{2\pi i}
\oint_{\sin\Gamma}
\frac{\cos((2k+1)\arcsin x)}{\sqrt{1-x^2}}
\frac{(\arcsin x)^{2p-1}}{x^{2n}}\,dx.
\end{cases}
\end{equation}

The factor
\[
\frac{\cos((2k+1)\arcsin x)}{\sqrt{1-x^2}}
\]
is exactly the Chebyshev factor. Indeed, if \(x=\sin\theta\), then
\[
U_{2k}(x)
=
U_{2k}(\sin\theta)
=
\frac{\sin((2k+1)(\pi/2-\theta))}{\cos\theta}
=
(-1)^k
\frac{\cos((2k+1)\theta)}{\cos\theta}.
\]
Thus
\[
\frac{\cos((2k+1)\arcsin x)}{\sqrt{1-x^2}}
=
(-1)^k U_{2k}(x).
\]
Therefore \eqref{eq:cd-after-sine-substitution-local} becomes
\[
\begin{cases}
\displaystyle
\sum_{m=p}^{n}
c_{2n-2m,n}
\frac{(2k+1)^{2m-2p}}{(2m-2p)!}
=
(-1)^{n+p+k}
\frac{1}{2\pi i}
\oint_{\sin\Gamma}
U_{2k}(x)(\arcsin x)^{2p}
\frac{dx}{x^{2n+1}},
\\[3.0em]
\displaystyle
\sum_{m=0}^{n-p}
d_{2m,n}
\frac{(2k+1)^{2n-2m-2p}}
{(2n-2m-2p)!}
=
(-1)^{n+p+k}
\frac{1}{2\pi i}
\oint_{\sin\Gamma}
U_{2k}(x)(\arcsin x)^{2p-1}
\frac{dx}{x^{2n}}.
\end{cases}
\]
By Cauchy's coefficient formula, this is exactly
\eqref{eq:c-sum-to-arcsine-chebyshev} and
\eqref{eq:d-sum-to-arcsine-chebyshev}.
\end{proof}

\begin{proposition}[Chebyshev--arcsine form of the shifted integrals]
\label[proposition]{prop:shifted_integrals_chebyshev_arcsine}
Let \(n\in\mathbb N\), \(0\le k<n\), and let \(U_m\) denote the Chebyshev
polynomial of the second kind, defined by
\[
U_m(\cos\theta)=\frac{\sin((m+1)\theta)}{\sin\theta}.
\]
Then
\begin{equation}\label{eq:shifted-zeta-chebyshev-arcsine}
\int_{0}^{\infty}
\frac{\sinh((2k+1)x)}{x\cosh^{2n+1}x}\,dx
=
\sum_{p=1}^{n}
(2^{2p+1}-1)
[z^{2n}]\,U_{2k}(z)(\arcsin z)^{2p}
\frac{\zeta(2p+1)}{\pi^{2p}},
\end{equation}
and
\begin{equation}\label{eq:shifted-beta-chebyshev-arcsine}
\int_{0}^{\infty}
\frac{\sinh((2k+1)x)}{x\cosh^{2n}x}\,dx
=
\sum_{p=1}^{n}
2^{2p}
[z^{2n-1}]\,U_{2k}(z)(\arcsin z)^{2p-1}
\frac{\beta(2p)}{\pi^{2p-1}}.
\end{equation}
\end{proposition}

\begin{proof}
We start from the contour evaluations 
\cite{talla_waffo_integral_2025}:
\begin{equation}\label{eq:talla-waffo-shifted-start}
\begin{cases}
\displaystyle
\int_{-\infty}^{\infty}
\frac{\sinh((2k+1)x)}{x\cosh^{2n+1}x}\,dx
=
\sum_{p=1}^{n}
\alpha_{p,k,n}
\frac{\zeta(2p+1)}{\pi^{2p}},
\\[2.4em]
\displaystyle
\int_{-\infty}^{\infty}
\frac{\sinh((2k+1)x)}{x\cosh^{2n}x}\,dx
=
\sum_{p=1}^{n}
f_{p,k,n}
\frac{\beta(2p)}{\pi^{2p-1}},
\end{cases}
\end{equation}
where
\begin{equation}\label{eq:talla-waffo-coefficients}
\begin{cases}
\displaystyle
\alpha_{p,k,n}
=
(-1)^{p+n+k}4^{p+1}
\left(1-\frac1{2^{2p+1}}\right)
\sum_{m=p}^{n}
c_{2n-2m,n}
\frac{(2k+1)^{2m-2p}}{(2m-2p)!},
\\[2.4em]
\displaystyle
f_{p,k,n}
=
(-1)^{n+k+p}2^{2p+1}
\sum_{m=0}^{n-p}
d_{2m,n}
\frac{(2k+1)^{2n-2m-2p}}
{(2n-2m-2p)!}.
\end{cases}
\end{equation}

By \cref{lemma:cd-sums-to-arcsine-chebyshev}, we have
\[
\sum_{m=p}^{n}
c_{2n-2m,n}
\frac{(2k+1)^{2m-2p}}{(2m-2p)!}
=
(-1)^{n+p+k}
[x^{2n}]\,U_{2k}(x)(\arcsin x)^{2p},
\]
and
\[
\sum_{m=0}^{n-p}
d_{2m,n}
\frac{(2k+1)^{2n-2m-2p}}
{(2n-2m-2p)!}
=
(-1)^{n+p+k}
[x^{2n-1}]\,U_{2k}(x)(\arcsin x)^{2p-1}.
\]
Substituting these two identities into \eqref{eq:talla-waffo-coefficients},
all signs cancel. Hence
\[
\alpha_{p,k,n}
=
4^{p+1}
\left(1-\frac1{2^{2p+1}}\right)
[x^{2n}]\,U_{2k}(x)(\arcsin x)^{2p}.
\]
Since
\[
4^{p+1}
\left(1-\frac1{2^{2p+1}}\right)
=
2(2^{2p+1}-1),
\]
we get
\[
\alpha_{p,k,n}
=
2(2^{2p+1}-1)
[x^{2n}]\,U_{2k}(x)(\arcsin x)^{2p}.
\]
Similarly,
\[
f_{p,k,n}
=
2^{2p+1}
[x^{2n-1}]\,U_{2k}(x)(\arcsin x)^{2p-1}.
\]

The integrands in \eqref{eq:talla-waffo-shifted-start} are even. Therefore
the integrals over \((0,\infty)\) are one half of the corresponding
integrals over \(\mathbb R\). Inserting the formulae for
\(\alpha_{p,k,n}\) and \(f_{p,k,n}\), and dividing by \(2\), proves
\eqref{eq:shifted-zeta-chebyshev-arcsine} and
\eqref{eq:shifted-beta-chebyshev-arcsine}.
\end{proof}

\begin{corollary}[Odd powers of \(\sinh x\)]
\label[corollary]{cor:odd-sinh-power-shifted-integrals}
Let \(n\in\mathbb N\) and \(0\le k<n\). Then
\begin{equation}
\label{eq:odd-sinh-zeta-general}
\int_0^\infty
\frac{\sinh^{2k+1}x}{x\cosh^{2n+1}x}\,dx
=
\sum_{p=1}^{n}
(2^{2p+1}-1)
[z^{2n}](z^2-1)^k(\arcsin z)^{2p}
\frac{\zeta(2p+1)}{\pi^{2p}},
\end{equation}
and
\begin{equation}
\label{eq:odd-sinh-beta-general}
\int_0^\infty
\frac{\sinh^{2k+1}x}{x\cosh^{2n}x}\,dx
=
\sum_{p=1}^{n}
2^{2p}[z^{2n-1}](z^2-1)^k(\arcsin z)^{2p-1}
\frac{\beta(2p)}{\pi^{2p-1}}.
\end{equation}
\end{corollary}

\begin{proof}
We first linearize the odd power of \(\sinh x\). From
\[
\sinh x=\frac{e^x-e^{-x}}2
\]
one obtains
\[
\sinh^{2k+1}x
=
\frac1{4^k}
\sum_{r=0}^{k}
(-1)^{k-r}
\binom{2k+1}{k-r}
\sinh((2r+1)x).
\]
Since \(0\le r\le k<n\), we may apply
\cref{prop:shifted_integrals_chebyshev_arcsine} to each term. Hence
\[
\int_0^\infty
\frac{\sinh^{2k+1}x}{x\cosh^{2n+1}x}\,dx
=
\sum_{p=1}^{n}
(2^{2p+1}-1)
[z^{2n}]
\left(
\frac1{4^k}
\sum_{r=0}^{k}
(-1)^{k-r}
\binom{2k+1}{k-r}
U_{2r}(z)
\right)
(\arcsin z)^{2p}
\frac{\zeta(2p+1)}{\pi^{2p}},
\]
and similarly
\[
\int_0^\infty
\frac{\sinh^{2k+1}x}{x\cosh^{2n}x}\,dx
=
\sum_{p=1}^{n}
2^{2p}
[z^{2n-1}]
\left(
\frac1{4^k}
\sum_{r=0}^{k}
(-1)^{k-r}
\binom{2k+1}{k-r}
U_{2r}(z)
\right)
(\arcsin z)^{2p-1}
\frac{\beta(2p)}{\pi^{2p-1}}.
\]

It remains to collapse the Chebyshev sum. Since
\[
\sinh((2r+1)x)=\sinh x\,U_{2r}(\cosh x),
\]
the same linearization gives, after division by \(\sinh x\),
\[
\sinh^{2k}x
=
\frac1{4^k}
\sum_{r=0}^{k}
(-1)^{k-r}
\binom{2k+1}{k-r}
U_{2r}(\cosh x).
\]
Putting \(z=\cosh x\), and using
\(\sinh^2x=\cosh^2x-1=z^2-1\), we get the polynomial identity
\[
\frac1{4^k}
\sum_{r=0}^{k}
(-1)^{k-r}
\binom{2k+1}{k-r}
U_{2r}(z)
=
(z^2-1)^k.
\]
Substituting this identity into the two preceding coefficient formulae proves
\eqref{eq:odd-sinh-zeta-general} and \eqref{eq:odd-sinh-beta-general}.
\end{proof}

\begin{corollary}[Integrals involving the reciprocal hyperbolic arctangent]
\label[corollary]{cor:reciprocal-arctanh-integrals}
Let \(n\in\mathbb N\). Then
\begin{equation}
\label{eq:reciprocal-arctanh-zeta}
\int_0^1
\frac{x^{2n-1}}{\operatorname{arctanh}x}\,dx
=
\sum_{p=1}^{n}
\left(2^{2p+1}-1\right)
[z^{2n}]
(z^2-1)^{n-1}(\arcsin z)^{2p}
\frac{\zeta(2p+1)}{\pi^{2p}},
\end{equation}
and
\begin{equation}
\label{eq:reciprocal-arctanh-beta}
\int_0^1
\frac{x^{2n-1}}
{\sqrt{1-x^2}\,\operatorname{arctanh}x}\,dx
=
\sum_{p=1}^{n}
2^{2p}
[z^{2n-1}]
(z^2-1)^{n-1}(\arcsin z)^{2p-1}
\frac{\beta(2p)}{\pi^{2p-1}}.
\end{equation}
\end{corollary}

\begin{proof}
We make the change of variables
\[
x=\tanh t,
\qquad
t=\operatorname{arctanh}x,
\qquad
dx=\frac{dt}{\cosh^2 t}.
\]
This substitution maps the interval \(0<x<1\) onto \(0<t<\infty\).

For the first integral, we obtain
\begin{align*}
\int_0^1
\frac{x^{2n-1}}{\operatorname{arctanh}x}\,dx
&=
\int_0^\infty
\frac{\tanh^{2n-1}t}{t\cosh^2t}\,dt
\\
&=
\int_0^\infty
\frac{\sinh^{2n-1}t}
{t\cosh^{2n+1}t}\,dt.
\end{align*}
Applying \eqref{eq:odd-sinh-zeta-general} with \(k=n-1\)
gives \eqref{eq:reciprocal-arctanh-zeta}.

Furthermore,
\[
\sqrt{1-\tanh^2t}
=
\frac{1}{\cosh t}.
\]
Consequently,
\begin{align*}
\int_0^1
\frac{x^{2n-1}}
{\sqrt{1-x^2}\,\operatorname{arctanh}x}\,dx
&=
\int_0^\infty
\frac{\tanh^{2n-1}t}
{t\cosh t}\,dt
\\
&=
\int_0^\infty
\frac{\sinh^{2n-1}t}
{t\cosh^{2n}t}\,dt.
\end{align*}
Applying \eqref{eq:odd-sinh-beta-general}, again with
\(k=n-1\), yields
\eqref{eq:reciprocal-arctanh-beta}.
\end{proof}

\section{Coefficient Formulae for $\displaystyle\int_{0}^{\infty}
\frac{\sinh(2kx)}{x\cosh^{n}x}\,dx$}\label[section]{sec:coefficient-formulae_double_angle}

\begin{lemma}[Reduction of binomial sums]
\label[lemma]{lem:chebyshev-binomial-collapse}
For every \(k\in\mathbb{N}\) with \(k\geq 1\), one has
\[
\sum_{j=0}^{k-1}
\binom{2k}{2j+1}
z^{2k-2j-1}(z^2-1)^j
=
U_{2k-1}(z),
\]
where \(U_{2k-1}\) denotes the Chebyshev polynomial of the
second kind.
\end{lemma}

\begin{proof}
Using the binomial expansion of \(\sinh(2kx)\), we obtain
\[
\sinh(2kx)
=
\sum_{j=0}^{k-1}
\binom{2k}{2j+1}
\sinh^{2j+1}(x)\cosh^{2k-2j-1}(x).
\]
Dividing by \(\sinh(x)\) gives
\[
\frac{\sinh(2kx)}{\sinh(x)}
=
\sum_{j=0}^{k-1}
\binom{2k}{2j+1}
\sinh^{2j}(x)\cosh^{2k-2j-1}(x).
\]
Since
\[
U_{2k-1}(\cosh x)
=
\frac{\sinh(2kx)}{\sinh x},
\]
and
\[
\sinh^2(x)=\cosh^2(x)-1,
\]
the substitution \(z=\cosh x\) yields
\[
U_{2k-1}(z)
=
\sum_{j=0}^{k-1}
\binom{2k}{2j+1}
z^{2k-2j-1}(z^2-1)^j.
\]
Both sides are polynomials in \(z\), so the identity holds
identically.
\end{proof}


\begin{proposition}
\label[proposition]{prop:sinh-over-cosh-integral}
Let \(k,m\in\mathbb{N}\) with \(k\geq 1\). Then
\[
\left\{
\begin{aligned}
\int_{0}^{\infty}
\frac{\sinh(2kx)}{x\cosh^{2m}x}\,dx
&=
\sum_{p=1}^{m-1}
\left(2^{2p+1}-1\right)\frac{\zeta(2p+1)}{\pi^{2p}}
[z^{2m-1}]
U_{2k-1}(z)(\arcsin z)^{2p}
,
&& m>k,
\\[1.5em]
\int_{0}^{\infty}
\frac{\sinh(2kx)}{x\cosh^{2m+1}x}\,dx
&=
\sum_{p=1}^{m}
2^{2p}\frac{\beta(2p)}{\pi^{2p-1}}
[z^{2m}]
U_{2k-1}(z)(\arcsin z)^{2p-1}
,
&& m\geq k.
\end{aligned}
\right.
\]
\end{proposition}

\begin{proof}
Set
\[
\left\{
\begin{aligned}
I_{k,m}^{\mathrm{e}}
&:=
\int_{0}^{\infty}
\frac{\sinh(2kx)}{x\cosh^{2m}x}\,dx,
\\[0.6em]
I_{k,m}^{\mathrm{o}}
&:=
\int_{0}^{\infty}
\frac{\sinh(2kx)}{x\cosh^{2m+1}x}\,dx.
\end{aligned}
\right.
\]

Near the origin, one has
\[
\left\{
\begin{aligned}
\frac{\sinh(2kx)}{x\cosh^{2m}x}
&=
2k+O(x^2),
\\
\frac{\sinh(2kx)}{x\cosh^{2m+1}x}
&=
2k+O(x^2),
\end{aligned}
\right.
\qquad x\to0.
\]
Moreover, as \(x\to\infty\),
\[
\left\{
\begin{aligned}
\frac{\sinh(2kx)}{x\cosh^{2m}x}
&=
O\left(
\frac{e^{-2(m-k)x}}{x}
\right),
\\[0.8em]
\frac{\sinh(2kx)}{x\cosh^{2m+1}x}
&=
O\left(
\frac{e^{-(2(m-k)+1)x}}{x}
\right).
\end{aligned}
\right.
\]
Consequently,
\[
\left\{
\begin{aligned}
I_{k,m}^{\mathrm{e}}
&\text{ converges absolutely for }m>k,
\\
I_{k,m}^{\mathrm{o}}
&\text{ converges absolutely for }m\geq k.
\end{aligned}
\right.
\]

Using the identity
\[
\sinh(2kx)
=
\sum_{j=0}^{k-1}
\binom{2k}{2j+1}
\sinh^{2j+1}(x)\cosh^{2k-2j-1}(x),
\]
we obtain
\[
\left\{
\begin{aligned}
I_{k,m}^{\mathrm{e}}
&=
\sum_{j=0}^{k-1}
\binom{2k}{2j+1}
\int_0^\infty
\frac{\sinh^{2j+1}(x)}
{x\cosh^{2(m-k+j)+1}(x)}
\,dx,
\\[1.2em]
I_{k,m}^{\mathrm{o}}
&=
\sum_{j=0}^{k-1}
\binom{2k}{2j+1}
\int_0^\infty
\frac{\sinh^{2j+1}(x)}
{x\cosh^{2(m-k+j+1)}(x)}
\,dx.
\end{aligned}
\right.
\]

For brevity, define
\[
a_p
:=
\left(2^{2p+1}-1\right)
\frac{\zeta(2p+1)}{\pi^{2p}},
\qquad
b_p
:=
2^{2p}
\frac{\beta(2p)}{\pi^{2p-1}}.
\]

Applying
\cref{cor:odd-sinh-power-shifted-integrals}
with the respective substitutions
\[
\left\{
\begin{aligned}
(r,\ell)
&=
(m-k+j,j),
\\
(r,\ell)
&=
(m-k+j+1,j),
\end{aligned}
\right.
\]
gives
\[
\left\{
\begin{aligned}
I_{k,m}^{\mathrm{e}}
&=
\sum_{j=0}^{k-1}
\binom{2k}{2j+1}
\sum_{p=1}^{m-k+j}
a_p
[z^{2(m-k+j)}]
\Bigl(
(z^2-1)^j(\arcsin z)^{2p}
\Bigr),
\\[1.2em]
I_{k,m}^{\mathrm{o}}
&=
\sum_{j=0}^{k-1}
\binom{2k}{2j+1}
\sum_{p=1}^{m-k+j+1}
b_p
[z^{2(m-k+j+1)-1}]
\Bigl(
(z^2-1)^j(\arcsin z)^{2p-1}
\Bigr).
\end{aligned}
\right.
\]

Since
\[
\operatorname{ord}_{z=0}(\arcsin z)^q=q,
\]
the relevant coefficients vanish whenever
\[
\left\{
\begin{aligned}
p&>m-k+j,
\\
p&>m-k+j+1,
\end{aligned}
\right.
\]
respectively. Hence, the inner sums may be extended to bounds
independent of \(j\).

Since
\[
\max_{0\leq j\leq k-1}(m-k+j)=m-1
\]
and
\[
\max_{0\leq j\leq k-1}(m-k+j+1)=m,
\]
interchanging the finite sums yields
\[
\left\{
\begin{aligned}
I_{k,m}^{\mathrm{e}}
&=
\sum_{p=1}^{m-1}
a_p
\sum_{j=0}^{k-1}
\binom{2k}{2j+1}
[z^{2(m-k+j)}]
\Bigl(
(z^2-1)^j(\arcsin z)^{2p}
\Bigr),
\\[1.2em]
I_{k,m}^{\mathrm{o}}
&=
\sum_{p=1}^{m}
b_p
\sum_{j=0}^{k-1}
\binom{2k}{2j+1}
[z^{2(m-k+j+1)-1}]
\Bigl(
(z^2-1)^j(\arcsin z)^{2p-1}
\Bigr).
\end{aligned}
\right.
\]

For every formal power series \(F(z)\), one has
\[
[z^r]F(z)
=
[z^{r+s}]
\bigl(z^sF(z)\bigr),
\qquad s\geq0.
\]
Taking
\[
s=2k-2j-1
\]
gives
\[
\left\{
\begin{aligned}
[z^{2(m-k+j)}]F(z)
&=
[z^{2m-1}]
\Bigl(
z^{2k-2j-1}F(z)
\Bigr),
\\[0.8em]
[z^{2(m-k+j+1)-1}]F(z)
&=
[z^{2m}]
\Bigl(
z^{2k-2j-1}F(z)
\Bigr).
\end{aligned}
\right.
\]

By linearity of the coefficient operator, it follows that
\[
\left\{
\begin{aligned}
I_{k,m}^{\mathrm{e}}
&=
\sum_{p=1}^{m-1}
a_p
[z^{2m-1}]
\Biggl(
(\arcsin z)^{2p}
\sum_{j=0}^{k-1}
\binom{2k}{2j+1}
z^{2k-2j-1}(z^2-1)^j
\Biggr),
\\[1.2em]
I_{k,m}^{\mathrm{o}}
&=
\sum_{p=1}^{m}
b_p
[z^{2m}]
\Biggl(
(\arcsin z)^{2p-1}
\sum_{j=0}^{k-1}
\binom{2k}{2j+1}
z^{2k-2j-1}(z^2-1)^j
\Biggr).
\end{aligned}
\right.
\]

By \cref{lem:chebyshev-binomial-collapse},
\[
\sum_{j=0}^{k-1}
\binom{2k}{2j+1}
z^{2k-2j-1}(z^2-1)^j
=
U_{2k-1}(z).
\]
Therefore,
\[
\left\{
\begin{aligned}
I_{k,m}^{\mathrm{e}}
&=
\sum_{p=1}^{m-1}
a_p
[z^{2m-1}]
\Bigl(
U_{2k-1}(z)(\arcsin z)^{2p}
\Bigr),
\\[1.2em]
I_{k,m}^{\mathrm{o}}
&=
\sum_{p=1}^{m}
b_p
[z^{2m}]
\Bigl(
U_{2k-1}(z)(\arcsin z)^{2p-1}
\Bigr).
\end{aligned}
\right.
\]

Substituting the definitions of \(a_p\) and \(b_p\), we finally obtain
\[
\left\{
\begin{aligned}
\int_{0}^{\infty}
\frac{\sinh(2kx)}{x\cosh^{2m}x}\,dx
&=
\sum_{p=1}^{m-1}
\left(2^{2p+1}-1\right)
[z^{2m-1}]
\Bigl(
U_{2k-1}(z)(\arcsin z)^{2p}
\Bigr)
\frac{\zeta(2p+1)}{\pi^{2p}},
\\[1.5em]
\int_{0}^{\infty}
\frac{\sinh(2kx)}{x\cosh^{2m+1}x}\,dx
&=
\sum_{p=1}^{m}
2^{2p}
[z^{2m}]
\Bigl(
U_{2k-1}(z)(\arcsin z)^{2p-1}
\Bigr)
\frac{\beta(2p)}{\pi^{2p-1}}.
\end{aligned}
\right.
\]
This completes the proof.
\end{proof}

\section{Coefficient Formulae for the logarithmic tangent integrals $\displaystyle\int_0^{\pi/4}
\frac{\sin(4nx)}{\ln(\tan x)}\,dx$ and $\displaystyle\int_0^{\pi/4}
\frac{\cos((4n-2)x)}{\ln(\tan x)}\,dx$}\label[section]{sec:coefficient-formulae-logarithmic-tangent-integrals}

\begin{lemma}[Tangent expansions of multiple angles]
\label[lemma]{lemma:tangent-multiple-angle-expansions}
Let \(n\in\mathbb N\). Then
\[
\cos(2nx)
=
\cos^{2n}x
\sum_{k=0}^{n}
(-1)^k
\binom{2n}{2k}
\tan^{2k}x
\]
and
\[
\sin(2nx)
=
\cos^{2n}x
\sum_{k=0}^{n-1}
(-1)^k
\binom{2n}{2k+1}
\tan^{2k+1}x .
\]
\end{lemma}

\begin{proof}
By de Moivre's formula,
\[
\cos(2nx)+i\sin(2nx)
=
(\cos x+i\sin x)^{2n}.
\]
Expanding the right-hand side gives
\[
(\cos x+i\sin x)^{2n}
=
\sum_{j=0}^{2n}
\binom{2n}{j}
i^j
\cos^{2n-j}x\,\sin^j x.
\]
Taking real parts, only even \(j=2k\) contribute. Hence
\[
\cos(2nx)
=
\sum_{k=0}^{n}
(-1)^k
\binom{2n}{2k}
\cos^{2n-2k}x\,\sin^{2k}x.
\]
Factoring out \(\cos^{2n}x\) yields
\[
\cos(2nx)
=
\cos^{2n}x
\sum_{k=0}^{n}
(-1)^k
\binom{2n}{2k}
\tan^{2k}x.
\]

Taking imaginary parts, only odd \(j=2k+1\) contribute. Therefore
\[
\sin(2nx)
=
\sum_{k=0}^{n-1}
(-1)^k
\binom{2n}{2k+1}
\cos^{2n-2k-1}x\,\sin^{2k+1}x.
\]
Again factoring out \(\cos^{2n}x\) gives
\[
\sin(2nx)
=
\cos^{2n}x
\sum_{k=0}^{n-1}
(-1)^k
\binom{2n}{2k+1}
\tan^{2k+1}x.
\]
\end{proof}

\begin{lemma}[Binomial sums of even Chebyshev polynomials]
\label[lemma]{lemma:even-chebyshev-binomial-sums}
Let \(n\ge 1\). Then the following identities hold in \(\mathbb Z[X]\):
\[
\sum_{k=0}^{n-1}
(-1)^k
\binom{4n}{2n-2k-1}
U_{2k}(X)
=
2^{2n-1}X^{2n-1}U_{2n-1}\!\left(\frac1X\right),
\]
and
\[
\sum_{k=0}^{n-1}
(-1)^k
\binom{4n-2}{2n-2k-2}
U_{2k}(X)
=
2^{2n-2}X^{2n-2}U_{2n-2}\!\left(\frac1X\right).
\]
Here \(U_m\) denotes the Chebyshev polynomial of the second kind.
\end{lemma}

\begin{proof}
We apply \cref{lemma:tangent-multiple-angle-expansions} with
\(\theta=\arctan x\), first with \(N=2n\) in the sine identity and then
with \(N=2n-1\) in the cosine identity. This gives
\[
\begin{cases}
\displaystyle
\sin(4n\arctan x)
=
\frac{1}{(1+x^2)^{2n}}
\sum_{r=0}^{2n-1}
(-1)^r\binom{4n}{2r+1}x^{2r+1},
\\[1.2em]
\displaystyle
\cos((4n-2)\arctan x)
=
\frac{1}{(1+x^2)^{2n-1}}
\sum_{r=0}^{2n-1}
(-1)^r\binom{4n-2}{2r}x^{2r}.
\end{cases}
\]
Pairing in both sums the terms \(r=n-k-1\) and \(r=n+k\), for
\(k=0,\ldots,n-1\), we obtain
\[
\begin{cases}
\displaystyle
\sin(4n\arctan x)
=
\frac{(-1)^{n-1}x^{2n}}{(1+x^2)^{2n}}
\sum_{k=0}^{n-1}
(-1)^k
\binom{4n}{2n-2k-1}
\left(x^{-2k-1}-x^{2k+1}\right),
\\[1.2em]
\displaystyle
\cos((4n-2)\arctan x)
=
\frac{(-1)^{n-1}x^{2n-1}}{(1+x^2)^{2n-1}}
\sum_{k=0}^{n-1}
(-1)^k
\binom{4n-2}{2n-2k-2}
\left(x^{-2k-1}-x^{2k+1}\right).
\end{cases}
\]

Now put \(x=e^u\) and \(X=\cosh u\). Since
\[
x^{-2k-1}-x^{2k+1}
=
-2\sinh((2k+1)u)
=
-2\sinh u\,U_{2k}(X),
\]
we get
\[
\begin{cases}
\displaystyle
\sin(4n\arctan e^u)
=
\frac{2(-1)^n\sinh u}{(2X)^{2n}}
\sum_{k=0}^{n-1}
(-1)^k
\binom{4n}{2n-2k-1}
U_{2k}(X),
\\[1.2em]
\displaystyle
\cos((4n-2)\arctan e^u)
=
\frac{2(-1)^n\sinh u}{(2X)^{2n-1}}
\sum_{k=0}^{n-1}
(-1)^k
\binom{4n-2}{2n-2k-2}
U_{2k}(X).
\end{cases}
\]

Let \(\alpha=2\arctan e^u\). Then
\[
\sin\alpha=\frac1X,
\qquad
\cos\alpha=-\frac{\sinh u}{X}.
\]
Using \(U_m(\cos t)=\sin((m+1)t)/\sin t\), with
\(t=\pi/2-\alpha\), gives
\[
\begin{cases}
\displaystyle
U_{2n-1}\!\left(\frac1X\right)
=
(-1)^n\frac{X\sin(4n\arctan e^u)}{\sinh u},
\\[1.2em]
\displaystyle
U_{2n-2}\!\left(\frac1X\right)
=
(-1)^n\frac{X\cos((4n-2)\arctan e^u)}{\sinh u}.
\end{cases}
\]
Substituting these into the preceding two identities yields the two claimed
formulae.

Finally, since \(X=\cosh u\) assumes infinitely many values and both sides
are polynomials in \(X\), the identities hold identically.
\end{proof}

\begin{proposition}[Compact coefficient form of the logarithmic tangent integrals]
\label[proposition]{proposition:logtan-integrals-compact}
Let \(U_m\) denote the Chebyshev polynomial of the second kind.  For every
\(n\in\mathbb N\),
\begin{equation}
\label{eq:sin-logtan-compact}
\int_0^{\pi/4}
\frac{\sin(4nx)}{\ln(\tan x)}\,dx
=
\frac{(-1)^n}{2}
\sum_{p=1}^{n}
(2^{2p+1}-1)
[x]\,
U_{2n-1}\!\left(\frac1x\right)(\arcsin x)^{2p}
\frac{\zeta(2p+1)}{\pi^{2p}},
\end{equation}
and
\begin{equation}
\label{eq:cos-logtan-compact}
\int_0^{\pi/4}
\frac{\cos((4n-2)x)}{\ln(\tan x)}\,dx
=
(-1)^n
\sum_{p=1}^{n}
2^{2p-1}
[x]\,
U_{2n-2}\!\left(\frac1x\right)(\arcsin x)^{2p-1}
\frac{\beta(2p)}{\pi^{2p-1}}.
\end{equation}
Here \([x]\) denotes the coefficient of \(x\) in the Laurent expansion at
\(x=0\).
\end{proposition}

\begin{proof}
We start from the logarithmic tangent evaluations
\cite{talla_waffo_integral_2025}:
\begin{equation}\label{eq:talla-waffo-logtan-start}
\begin{cases}
\displaystyle
\int_0^{\pi/4}
\frac{\sin(4nx)}{\ln(\tan x)}\,dx
=
\sum_{p=1}^{n}
q_{p,n}\frac{\zeta(2p+1)}{\pi^{2p}},
\\[2.4em]
\displaystyle
\int_0^{\pi/4}
\frac{\cos((4n-2)x)}{\ln(\tan x)}\,dx
=
\sum_{p=1}^{n}
\eta_{p,n}\frac{\beta(2p)}{\pi^{2p-1}},
\end{cases}
\end{equation}
where
\begin{equation}\label{eq:talla-waffo-logtan-coefficients}
\begin{cases}
\displaystyle
q_{p,n}
=
\frac{1}{4^n}
\sum_{k=0}^{n-1}
\binom{4n}{2n-2k-1}
\sum_{m=p}^{n}
(-1)^p c_{2n-2m,n}
\frac{(2^{2p+1}-1)(2k+1)^{2m-2p}}{(2m-2p)!},
\\[3.0em]
\displaystyle
\eta_{p,n}
=
\frac{(-1)^p}{2^{2n-2p-1}}
\sum_{m=0}^{n-p}
\frac{d_{2m,n}}{(2n-2m-2p)!}
\sum_{k=0}^{n-1}
\binom{4n-2}{2n-2k-2}
(2k+1)^{2n-2m-2p}.
\end{cases}
\end{equation}

By \cref{lemma:cd-sums-to-arcsine-chebyshev}, this gives
\begin{equation}\label{eq:logtan-coefficients-arcsine-step}
\begin{cases}
\displaystyle
q_{p,n}
=
(2^{2p+1}-1)
[x^{2n}]
\left\{
\frac{(-1)^n}{4^n}
\sum_{k=0}^{n-1}
(-1)^k
\binom{4n}{2n-2k-1}
U_{2k}(x)
\right\}
(\arcsin x)^{2p},
\\[3.0em]
\displaystyle
\eta_{p,n}
=
2^{2p}
[x^{2n-1}]
\left\{
\frac{(-1)^n}{2^{2n-1}}
\sum_{k=0}^{n-1}
(-1)^k
\binom{4n-2}{2n-2k-2}
U_{2k}(x)
\right\}
(\arcsin x)^{2p-1}.
\end{cases}
\end{equation}

Using \cref{lemma:even-chebyshev-binomial-sums}, we obtain
\begin{equation}\label{eq:logtan-coefficients-collapsed-step}
\begin{cases}
\displaystyle
q_{p,n}
=
\frac{(-1)^n}{2}
(2^{2p+1}-1)
[x^{2n}]
x^{2n-1}
U_{2n-1}\!\left(\frac1x\right)
(\arcsin x)^{2p},
\\[3.0em]
\displaystyle
\eta_{p,n}
=
(-1)^n2^{2p-1}
[x^{2n-1}]
x^{2n-2}
U_{2n-2}\!\left(\frac1x\right)
(\arcsin x)^{2p-1}.
\end{cases}
\end{equation}

For every Laurent series \(F(x)\) and every integer \(r\), we use the
coefficient shift
\[
[x^r]\,x^{r-1}F(x)=[x]\,F(x).
\]
Hence
\begin{equation}\label{eq:logtan-coefficients-compact-step}
\begin{cases}
\displaystyle
q_{p,n}
=
\frac{(-1)^n}{2}
(2^{2p+1}-1)
[x]\,
U_{2n-1}\!\left(\frac1x\right)
(\arcsin x)^{2p},
\\[3.0em]
\displaystyle
\eta_{p,n}
=
(-1)^n2^{2p-1}
[x]\,
U_{2n-2}\!\left(\frac1x\right)
(\arcsin x)^{2p-1}.
\end{cases}
\end{equation}

Substituting \eqref{eq:logtan-coefficients-compact-step} into
\eqref{eq:talla-waffo-logtan-start} proves
\eqref{eq:sin-logtan-compact} and \eqref{eq:cos-logtan-compact}.
\end{proof}

\begin{remark}[The first beta coefficient]
\label[remark]{rem:first-beta-coefficient}
The first three cases of \eqref{eq:cos-logtan-compact} are
\begin{align}
\int_0^{\pi/4}
\frac{\cos(2x)}{\ln(\tan x)}\,dx
&=
-\frac{2\beta(2)}{\pi},
\label{eq:cos2x-logtan}
\\
\int_0^{\pi/4}
\frac{\cos(6x)}{\ln(\tan x)}\,dx
&=
-\frac{2\beta(2)}{3\pi}
+
\frac{32\beta(4)}{\pi^3},
\label{eq:cos6x-logtan}
\\
\int_0^{\pi/4}
\frac{\cos(10x)}{\ln(\tan x)}\,dx
&=
-\frac{2\beta(2)}{5\pi}
+
\frac{32\beta(4)}{\pi^3}
-
\frac{512\beta(6)}{\pi^5}.
\label{eq:cos10x-logtan}
\end{align}
All three evaluations already appear in
\cite{talla_waffo_integral_2025}.

The first three examples exhibit a particularly simple pattern for the
coefficient of \(\beta(2)/\pi\). Indeed, the term corresponding to
\(p=1\) in \eqref{eq:cos-logtan-compact} is
\[
(-1)^n
2[x]\,
U_{2n-2}\!\left(\frac1x\right)\arcsin x
\frac{\beta(2)}{\pi},
\]
and one has
\begin{equation}
\label{eq:first-chebyshev-arcsin-coefficient}
[x]\,
U_{2n-2}\!\left(\frac1x\right)\arcsin x
=
\frac{(-1)^{n-1}}{2n-1}.
\end{equation}
Consequently, the coefficient of \(\beta(2)/\pi\) is
\[
-\frac{2}{2n-1}.
\]

A proof of this coefficient is already given
in \cite{talla_waffo_integral_2025}. In
\cref{prop:first-chebyshev-arcsin-coefficient}, we provide a different
proof based directly on the new coefficient-extraction formula
\eqref{eq:cos-logtan-compact}.
\end{remark}


\begin{proposition}[The first Chebyshev--arcsine coefficient]
\label[proposition]{prop:first-chebyshev-arcsin-coefficient}
For every integer \(n\geq1\),
\begin{equation}
\label{eq:first-chebyshev-arcsin-coefficient}
\boxed{
[x]\,
U_{2n-2}\!\left(\frac1x\right)\arcsin x
=
\frac{(-1)^{n-1}}{2n-1}
}.
\end{equation}
\end{proposition}

\begin{proof} Let \(C_\rho=\{|z|=\rho\}\) be a sufficiently small positively oriented circle around the origin. By Cauchy's coefficient formula for Laurent series, \[ [x]\, U_{2n-2}\!\left(\frac1x\right)\arcsin x = \frac{1}{2\pi i} \oint_{C_\rho} \frac{\arcsin z}{z^2} U_{2n-2}\!\left(\frac1z\right)\,dz. \] We make the local biholomorphic substitution \[ \frac1z=\frac{w+w^{-1}}2, \qquad\text{equivalently}\qquad z=\frac{2w}{1+w^2}. \] Then \[ dz=\frac{2(1-w^2)}{(1+w^2)^2}\,dw, \qquad \frac{dz}{z^2} = \frac{1-w^2}{2w^2}\,dw. \] Hence \[ [x]\, U_{2n-2}\!\left(\frac1x\right)\arcsin x = \frac{1}{2\pi i} \oint_{\Gamma_\rho} \arcsin\!\left(\frac{2w}{1+w^2}\right) U_{2n-2}\!\left(\frac{w+w^{-1}}2\right) \frac{1-w^2}{2w^2}\,dw, \] where \(\Gamma_\rho\) is a sufficiently small positively oriented contour around the origin. Using \[ \frac{w+w^{-1}}2=\cosh(\ln w) \quad \text{and} \quad U_m(\cosh \xi) = \frac{\sinh((m+1)\xi)}{\sinh \xi}, \] we obtain \[ \begin{aligned} U_{2n-2}\!\left(\frac{w+w^{-1}}2\right) &= U_{2n-2}\!\bigl(\cosh(\ln w)\bigr)\\ &= \frac{\sinh\!\bigl((2n-1)\ln w\bigr)} {\sinh(\ln w)}\\ &= \frac{e^{(2n-1)\ln w}-e^{-(2n-1)\ln w}} {e^{\ln w}-e^{-\ln w}}\\ &= \frac{w^{2n-1}-w^{-(2n-1)}}{w-w^{-1}}. \end{aligned} \] Therefore, \[ \begin{aligned} [x]\, U_{2n-2}\!\left(\frac1x\right)\arcsin x &= \frac{1}{2\pi i} \oint_{\Gamma_\rho} \arcsin\!\left(\frac{2w}{1+w^2}\right) \frac{w^{2n-1}-w^{-(2n-1)}}{w-w^{-1}} \frac{1-w^2}{2w^2}\,dw\\ &= \frac{1}{2\pi i} \oint_{\Gamma_\rho} \frac12\left(w^{-2n}-w^{2n-2}\right) \arcsin\!\left(\frac{2w}{1+w^2}\right)\,dw. \end{aligned} \] Since, as an identity of holomorphic germs at the origin, \[ \arcsin\!\left(\frac{2w}{1+w^2}\right) = 2\arctan w, \] it follows that \[ [x]\, U_{2n-2}\!\left(\frac1x\right)\arcsin x = \frac{1}{2\pi i} \oint_{\Gamma_\rho} \left(w^{-2n}-w^{2n-2}\right)\arctan w\,dw. \] Now \[ \arctan w = \sum_{k\geq0} \frac{(-1)^k}{2k+1}w^{2k+1}. \] The term \(w^{2n-2}\arctan w\) is holomorphic at the origin, while \[ \operatorname*{Res}_{w=0} \left(w^{-2n}\arctan w\right) = [w^{2n-1}]\arctan w = \frac{(-1)^{n-1}}{2n-1}. \] Consequently, \[ [x]\, U_{2n-2}\!\left(\frac1x\right)\arcsin x = \frac{(-1)^{n-1}}{2n-1}. \] \end{proof}

\begin{remark}[A parity-free formulation]
\label[remark]{rem:parity-free-terminal-coefficient}
The terminal coefficients in the sine and cosine expansions were already
determined in \cite{talla_waffo_integral_2025}. More precisely, one has
\begin{equation}
\label{eq:previous-terminal-coefficients}
q_{n,n}
=
(-1)^n
\bigl(2^{2n+1}-1\bigr)4^{n-1},
\qquad
\eta_{n,n}
=
(-1)^n2^{4n-3}.
\end{equation}

We now compare these formulae with the coefficient-extraction identities
\eqref{eq:sin-logtan-compact} and \eqref{eq:cos-logtan-compact}.

The coefficient of
\[
\frac{\zeta(2n+1)}{\pi^{2n}}
\]
in \eqref{eq:sin-logtan-compact}, obtained by taking \(p=n\), is
\[
\frac{(-1)^n}{2}
\bigl(2^{2n+1}-1\bigr)
[x]\,
U_{2n-1}\!\left(\frac1x\right)
(\arcsin x)^{2n}.
\]
Comparing this expression with \(q_{n,n}\) in
\eqref{eq:previous-terminal-coefficients} gives
\begin{equation}
\label{eq:terminal-sine-coefficient-extraction}
[x]\,
U_{2n-1}\!\left(\frac1x\right)
(\arcsin x)^{2n}
=
2^{2n-1}.
\end{equation}

Similarly, the coefficient of
\[
\frac{\beta(2n)}{\pi^{2n-1}}
\]
in \eqref{eq:cos-logtan-compact} is
\[
(-1)^n2^{2n-1}
[x]\,
U_{2n-2}\!\left(\frac1x\right)
(\arcsin x)^{2n-1}.
\]
Comparison with \(\eta_{n,n}\) in
\eqref{eq:previous-terminal-coefficients} yields
\begin{equation}
\label{eq:terminal-cosine-coefficient-extraction}
[x]\,
U_{2n-2}\!\left(\frac1x\right)
(\arcsin x)^{2n-1}
=
2^{2n-2}.
\end{equation}

The two identities
\eqref{eq:terminal-sine-coefficient-extraction} and
\eqref{eq:terminal-cosine-coefficient-extraction} admit a single
formulation that is independent of parity. Namely, setting \(r=2n\) in
\eqref{eq:terminal-sine-coefficient-extraction} and \(r=2n-1\) in
\eqref{eq:terminal-cosine-coefficient-extraction}, respectively, gives
\begin{equation}
\label{eq:parity-free-terminal-coefficient}
[x]\,
U_{r-1}\!\left(\frac1x\right)
(\arcsin x)^r
=
2^{r-1},
\qquad r\geq1.
\end{equation}

Thus the terminal coefficients in both
\eqref{eq:sin-logtan-compact} and
\eqref{eq:cos-logtan-compact} are governed by the same
parity-independent coefficient-extraction identity.

Although \eqref{eq:parity-free-terminal-coefficient} follows here from
the terminal coefficient formulae established in
\cite{talla_waffo_integral_2025}, in
\cref{prop:parity-free-terminal-coefficient} we give a new direct proof
based solely on the coefficient-extraction representation. This single
argument proves simultaneously both
\eqref{eq:terminal-sine-coefficient-extraction} and
\eqref{eq:terminal-cosine-coefficient-extraction}.
\end{remark}


\begin{proposition}[The parity-free terminal coefficient]
\label[proposition]{prop:parity-free-terminal-coefficient}
For every integer \(r\geq1\),
\begin{equation}
\label{eq:parity-free-terminal-coefficient-proposition}
\boxed{
[x]\,
U_{r-1}\!\left(\frac1x\right)
(\arcsin x)^r
=
2^{r-1}
}.
\end{equation}
\end{proposition}

\begin{proof}
The Chebyshev polynomial \(U_{r-1}\) has degree \(r-1\), leading
coefficient \(2^{r-1}\), and contains only powers of the same parity as
\(r-1\). Therefore, as a Laurent expansion at \(x=0\),
\[
U_{r-1}\!\left(\frac1x\right)
=
2^{r-1}x^{1-r}+O(x^{3-r}).
\]

Moreover,
\[
\arcsin x=x+O(x^3),
\]
and hence
\[
(\arcsin x)^r
=
x^r+O(x^{r+2}).
\]
Multiplying these expansions gives
\[
\begin{aligned}
U_{r-1}\!\left(\frac1x\right)(\arcsin x)^r
&=
\left(2^{r-1}x^{1-r}+O(x^{3-r})\right)
\left(x^r+O(x^{r+2})\right)\\
&=
2^{r-1}x+O(x^3).
\end{aligned}
\]
In particular, no other term can contribute to the coefficient of \(x\).
Therefore
\[
[x]\,
U_{r-1}\!\left(\frac1x\right)(\arcsin x)^r
=
2^{r-1},
\]
which proves
\eqref{eq:parity-free-terminal-coefficient-proposition}.
\end{proof}

\section{The general tangent power integral on the half-line}\label[section]{section:general-tangent-integral}

Throughout this section, let \(m,n\geq1\), \(m\geq n\), and assume that
\(m+n\) is even. Put
\[
L:=(1-y^2)\frac{d}{dy},
\qquad
j_-:=\frac{m-n}{2},
\qquad
j_+:=\frac{m+n-2}{2}.
\]
Notice that \(j_-\) and \(j_+\) are nonnegative integers.

\begin{lemma}[The auxiliary polynomial]
\label[lemma]{lem:auxiliary-polynomial-general-tanh}
There is a unique polynomial
\[
\Phi_{m,n}(z)
=
\sum_{j=j_-}^{j_+}a_{m,n,j}z^j
\]
such that
\begin{equation}
\label{eq:Phi-mn-definition}
L^n y^{m+1}
=
y(1-y^2)\Phi_{m,n}(y^2).
\end{equation}
The endpoint coefficients are nonzero and satisfy
\begin{equation}
\label{eq:Phi-mn-endpoint-coefficients}
a_{m,n,j_-}
=
\frac{(m+1)!}{(m+1-n)!},
\qquad
a_{m,n,j_+}
=
(-1)^{n+1}\frac{(m+n)!}{m!}.
\end{equation}
Equivalently,
\begin{equation}
\label{eq:Dn-tanh-mn-Phi}
\frac{d^n}{dx^n}\tanh^{m+1}x
=
\tanh x\,\sech^2x\,
\Phi_{m,n}(\tanh^2x).
\end{equation}
\end{lemma}

\begin{proof}
For every integer \(r\geq1\), one has
\[
Ly^r
=
r y^{r-1}-r y^{r+1}.
\]
Consequently, every monomial that can occur in \(L^n y^{m+1}\) has
exponent belonging to
\[
\{m+1-n,m+3-n,\ldots,m+n+1\}.
\]

The smallest exponent \(m+1-n\) is obtained by choosing the
degree-lowering term at every application of \(L\). Its coefficient is
\[
(m+1)m\cdots(m+2-n)
=
\frac{(m+1)!}{(m+1-n)!},
\]
and is therefore nonzero.

Similarly, the largest exponent \(m+n+1\) is obtained by choosing the
degree-raising term at every application of \(L\). Its coefficient is
\[
(-1)^n(m+1)(m+2)\cdots(m+n)
=
(-1)^n\frac{(m+n)!}{m!},
\]
and is also nonzero.

The operator \(L\) changes parity, since \(1-y^2\) is even and
\(d/dy\) changes parity. Thus \(L^n y^{m+1}\) has parity \(m+1-n\).
Since \(m+n\) is even, \(m+1-n\) is odd. Hence \(L^n y^{m+1}\) is an
odd polynomial.

Moreover, since \(n\geq1\), the final application of \(L\) contributes
the factor \(1-y^2\). Thus
\[
L^n y^{m+1}
=
(1-y^2)R_{m,n}(y)
\]
for some polynomial \(R_{m,n}\). Since \(L^n y^{m+1}\) is odd and
\(1-y^2\) is even, \(R_{m,n}\) is odd. Therefore there is a unique
polynomial \(\Phi_{m,n}\) such that
\[
R_{m,n}(y)
=
y\Phi_{m,n}(y^2),
\]
which proves \eqref{eq:Phi-mn-definition}.

Let \(j_{\min}\) and \(j_{\max}\) denote respectively the smallest and
largest indices occurring in \(\Phi_{m,n}\) with nonzero coefficient.
Then
\[
y(1-y^2)\Phi_{m,n}(y^2)
=
\sum_j a_{m,n,j}
\bigl(y^{2j+1}-y^{2j+3}\bigr).
\]
Its smallest exponent is \(2j_{\min}+1\), while its largest exponent is
\(2j_{\max}+3\). Comparing these with the extreme exponents found above
gives
\[
2j_{\min}+1=m+1-n,
\qquad
2j_{\max}+3=m+n+1.
\]
Hence
\[
j_{\min}
=
\frac{m-n}{2}
=
j_-,
\qquad
j_{\max}
=
\frac{m+n-2}{2}
=
j_+.
\]

Comparing the coefficients of the smallest and largest powers of \(y\)
also gives
\[
a_{m,n,j_-}
=
\frac{(m+1)!}{(m+1-n)!}
\]
and
\[
-a_{m,n,j_+}
=
(-1)^n\frac{(m+n)!}{m!},
\]
which proves \eqref{eq:Phi-mn-endpoint-coefficients}.

Finally, if \(y=\tanh x\), then
\[
\frac{dy}{dx}
=
1-y^2
\]
and therefore
\[
\frac{d}{dx}
=
(1-y^2)\frac{d}{dy}
=
L.
\]
Substituting \(y=\tanh x\) in \eqref{eq:Phi-mn-definition} gives
\eqref{eq:Dn-tanh-mn-Phi}.
\end{proof}


\begin{lemma}[Repeated integration by parts]
\label[lemma]{lem:general-tanh-half-line-ibp}
One has
\begin{equation}
\label{eq:general-tanh-half-line-ibp}
\int_0^\infty
\frac{\tanh^{m+1}x}{x^{n+1}}\,dx
=
\frac1{n!}
\int_0^\infty
\frac{D^n(\tanh^{m+1}x)}{x}\,dx.
\end{equation}
Writing
\[
\Phi_{m,n}(z)
=
\sum_{j=j_-}^{j_+}a_{m,n,j}z^j,
\]
one has
\begin{equation}
\label{eq:general-tanh-Phi-reduction}
\int_0^\infty
\frac{\tanh^{m+1}x}{x^{n+1}}\,dx
=
\frac1{n!}
\sum_{j=j_-}^{j_+}
a_{m,n,j}
\int_0^\infty
\frac{\sinh^{2j+1}x}
{x\cosh^{2j+3}x}\,dx.
\end{equation}
\end{lemma}

\begin{proof}
Since
\[
D^n\!\left(\frac1x\right)
=
(-1)^n\frac{n!}{x^{n+1}},
\]
we integrate by parts \(n\) times on \((\varepsilon,R)\) and then let
\(\varepsilon\to0^+\) and \(R\to\infty\).

The boundary terms vanish. Near \(0\),
\[
D^r(\tanh^{m+1}x)
=
O(x^{m+1-r})
\qquad
(0\leq r\leq n).
\]
Every boundary term at \(0\) is therefore
\[
O(x^{m+1-n}),
\]
and hence tends to zero because \(m\geq n\).

At \(\infty\), the boundary term involving \(\tanh^{m+1}x\) tends to
zero because it is multiplied by a negative power of \(x\). All boundary
terms involving at least one derivative of \(\tanh^{m+1}x\) tend to zero
exponentially. This proves
\eqref{eq:general-tanh-half-line-ibp}.

Using \cref{lem:auxiliary-polynomial-general-tanh}, we obtain
\[
D^n(\tanh^{m+1}x)
=
\tanh x\,\sech^2x\,
\Phi_{m,n}(\tanh^2x).
\]
Expanding \(\Phi_{m,n}\) gives
\[
D^n(\tanh^{m+1}x)
=
\sum_{j=j_-}^{j_+}
a_{m,n,j}
\tanh^{2j+1}x\,\sech^2x.
\]
Since
\[
\tanh^{2j+1}x\,\sech^2x
=
\frac{\sinh^{2j+1}x}{\cosh^{2j+3}x},
\]
equation \eqref{eq:general-tanh-Phi-reduction} follows.
\end{proof}


\begin{lemma}[General cotangent collapse]
\label[lemma]{lem:general-cotangent-collapse}
Let \(p\geq1\), and define
\begin{equation}
\label{eq:B-mnp-definition}
B_{m,n,p}
:=
\sum_{j=j_-}^{j_+}
a_{m,n,j}
[z^{2j+2}]
(z^2-1)^j(\arcsin z)^{2p}.
\end{equation}
Then \(B_{m,n,p}=0\) for \(2p<n\), and, for \(2p\geq n\),
\begin{equation}
\label{eq:B-mnp-collapsed}
B_{m,n,p}
=
(-1)^{(m-n)/2}
\frac{(2p)!}{(2p-n)!}
[u^{m+n-2p}](u\cot u)^{m+1}.
\end{equation}
\end{lemma}

\begin{proof} We first rewrite \(B_{m,n,p}\) as a single Laurent coefficient. Since \[ \Phi_{m,n}(w)=\sum_{j=j_-}^{j_+}a_{m,n,j}w^j, \] we have \[ \Phi_{m,n}\!\left(1-\frac1{z^2}\right) = \sum_{j=j_-}^{j_+} a_{m,n,j} \left(1-\frac1{z^2}\right)^j. \] Moreover, since \[ (z^2-1)^j=z^{2j}\left(1-\frac1{z^2}\right)^j, \] it follows that \[ [z^{2j+2}](z^2-1)^j(\arcsin z)^{2p} = [z^2]\left(1-\frac1{z^2}\right)^j(\arcsin z)^{2p}. \] Hence \begin{equation} \label{eq:B-mnp-Laurent} B_{m,n,p} = [z^2](\arcsin z)^{2p} \Phi_{m,n}\!\left(1-\frac1{z^2}\right), \end{equation} where the coefficient is understood in the Laurent expansion at \(z=0\). Let \(C_\rho=\{|z|=\rho\}\) be positively oriented and sufficiently small that \(\arcsin z\) is holomorphic on and inside \(C_\rho\). By Cauchy's coefficient formula for Laurent series, \begin{equation} \label{eq:B-mnp-contour-z} B_{m,n,p} = \frac1{2\pi i} \oint_{C_\rho} \frac{(\arcsin z)^{2p}}{z^3} \Phi_{m,n}\!\left(1-\frac1{z^2}\right)\,dz, \end{equation} since division by \(z^3\) extracts the coefficient of \(z^2\). We now make the local biholomorphic change of variables \(z=\sin u\) near \(u=0\). If \(\widetilde C_\rho\) is the positively oriented preimage of \(C_\rho\), then \[ \arcsin z=u,\qquad dz=\cos u\,du,\qquad 1-\frac1{z^2}=1-\csc^2u=-\cot^2u. \] Therefore \eqref{eq:B-mnp-contour-z} becomes \[ B_{m,n,p} = \frac1{2\pi i} \oint_{\widetilde C_\rho} u^{2p}\frac{\cos u}{\sin^3u} \Phi_{m,n}(-\cot^2u)\,du. \] Since \(\cos u/\sin^3u=\cot u\,\csc^2u\), this gives \[ B_{m,n,p} = \frac1{2\pi i} \oint_{\widetilde C_\rho} u^{2p}\cot u\,\csc^2u\, \Phi_{m,n}(-\cot^2u)\,du. \] The integrand is meromorphic near \(u=0\) and has no singularities in a sufficiently small punctured neighbourhood of \(0\). Hence the contour may be deformed to any sufficiently small positively oriented circle \(C_r=\{|u|=r\}\), yielding \begin{equation} \label{eq:B-mnp-contour-u} B_{m,n,p} = \frac1{2\pi i} \oint_{C_r} u^{2p}\cot u\,\csc^2u\, \Phi_{m,n}(-\cot^2u)\,du. \end{equation} We next identify the cotangent derivative. Starting from \[ L^ny^{m+1} = y(1-y^2)\Phi_{m,n}(y^2), \qquad L=(1-y^2)\frac{d}{dy}, \] put \(y=iv\). Since \(1-y^2=1+v^2\) and \(\frac{d}{dy}=\frac1i\frac{d}{dv}\), one has \[ L=\frac1i(1+v^2)\frac{d}{dv}. \] Thus \[ \frac1{i^n} \left((1+v^2)\frac{d}{dv}\right)^n (iv)^{m+1} = iv(1+v^2)\Phi_{m,n}(-v^2). \] Simplifying the powers of \(i\), and using that \(m+n\) is even, gives \begin{equation} \label{eq:Phi-mn-v-form} \left((1+v^2)\frac{d}{dv}\right)^n v^{m+1} = (-1)^{(m-n)/2} v(1+v^2)\Phi_{m,n}(-v^2). \end{equation} Now put \(v=\cot u\). Since \[ \frac{dv}{du}=-\csc^2u=-(1+v^2), \] we have \(\frac{d}{du}=-(1+v^2)\frac{d}{dv}\). Applying this identity \(n\) times to \eqref{eq:Phi-mn-v-form} yields \begin{equation} \label{eq:cot-derivative-general-mn} \frac{d^n}{du^n}\cot^{m+1}u = (-1)^{n+(m-n)/2} \cot u\,\csc^2u\, \Phi_{m,n}(-\cot^2u). \end{equation} Substituting \eqref{eq:cot-derivative-general-mn} into \eqref{eq:B-mnp-contour-u}, we obtain \begin{equation} \label{eq:B-mnp-contour-derivative} B_{m,n,p} = \frac{(-1)^{n+(m-n)/2}}{2\pi i} \oint_{C_r} u^{2p} \frac{d^n}{du^n}\cot^{m+1}u\,du. \end{equation} We now integrate by parts \(n\) times along \(C_r\). For meromorphic functions \(F\) and \(G\) on a neighbourhood of \(C_r\), \[ \oint_{C_r}F(u)G'(u)\,du = -\oint_{C_r}F'(u)G(u)\,du, \] because the integral of \((FG)'\) over a closed contour is zero. Iterating, \[ \oint_{C_r}F(u)G^{(n)}(u)\,du = (-1)^n \oint_{C_r}F^{(n)}(u)G(u)\,du. \] Taking \(F(u)=u^{2p}\) and \(G(u)=\cot^{m+1}u\), \eqref{eq:B-mnp-contour-derivative} becomes \begin{equation} \label{eq:B-mnp-after-parts} B_{m,n,p} = \frac{(-1)^{(m-n)/2}}{2\pi i} \oint_{C_r} \frac{d^n}{du^n}(u^{2p}) \cot^{m+1}u\,du. \end{equation} If \(2p<n\), then \(\frac{d^n}{du^n}u^{2p}=0\), and hence \(B_{m,n,p}=0\). If \(2p\geq n\), then \[ \frac{d^n}{du^n}u^{2p} = \frac{(2p)!}{(2p-n)!}u^{2p-n}, \] so \begin{equation} \label{eq:B-mnp-before-final-coefficient} B_{m,n,p} = (-1)^{(m-n)/2} \frac{(2p)!}{(2p-n)!} \frac1{2\pi i} \oint_{C_r} u^{2p-n}\cot^{m+1}u\,du. \end{equation} Finally, \[ u^{2p-n}\cot^{m+1}u = u^{2p-n-m-1}(u\cot u)^{m+1}. \] Since \(u\cot u\) is holomorphic at \(u=0\) and has value \(1\), Cauchy's coefficient formula gives \[ \frac1{2\pi i} \oint_{C_r} u^{2p-n}\cot^{m+1}u\,du = [u^{m+n-2p}](u\cot u)^{m+1}. \] Substitution into \eqref{eq:B-mnp-before-final-coefficient} proves \eqref{eq:B-mnp-collapsed}. \end{proof} 

\begin{proposition}
[General zeta formula for the hyperbolic tangent power integral]
\label[proposition]{prop:general-tangent-power-integral}
Let \(m,n\geq1\), \(m\geq n\), and suppose that \(m+n\) is even. Then
\begin{equation}
\label{eq:general-tangent-power-integral-formula}
\boxed{
\int_0^\infty
\frac{\tanh^{m+1}x}{x^{n+1}}\,dx
=
(-1)^{(m-n)/2}
\sum_{p=\lceil n/2\rceil}^{(m+n)/2}
\binom{2p}{n}
(2^{2p+1}-1)
\frac{\zeta(2p+1)}{\pi^{2p}}
[u^{m+n-2p}](u\cot u)^{m+1}
}
\end{equation}
\end{proposition}

\begin{proof} By \cref{lem:general-tanh-half-line-ibp}, \[ \int_0^\infty \frac{\tanh^{m+1}x}{x^{n+1}}\,dx = \frac1{n!} \sum_{j=j_-}^{j_+} a_{m,n,j} \int_0^\infty \frac{\sinh^{2j+1}x}{x\cosh^{2j+3}x}\,dx. \] Applying \eqref{eq:odd-sinh-zeta-general} with \(k=j\) and \(n=j+1\), we obtain \[ \int_0^\infty \frac{\tanh^{m+1}x}{x^{n+1}}\,dx = \frac1{n!} \sum_{p\geq1} (2^{2p+1}-1) \frac{\zeta(2p+1)}{\pi^{2p}} B_{m,n,p}. \] By \cref{lem:general-cotangent-collapse}, one has \(B_{m,n,p}=0\) whenever \(2p<n\), while for \(2p\geq n\), \[ B_{m,n,p} = (-1)^{(m-n)/2} \frac{(2p)!}{(2p-n)!} [u^{m+n-2p}](u\cot u)^{m+1}. \] Hence \(p\geq\lceil n/2\rceil\). Moreover, the coefficient extraction vanishes when \(m+n-2p<0\), so only \(p\leq(m+n)/2\) contributes. Using \[ \frac1{n!}\frac{(2p)!}{(2p-n)!} = \binom{2p}{n}, \] we conclude that \[ \int_0^\infty \frac{\tanh^{m+1}x}{x^{n+1}}\,dx = (-1)^{(m-n)/2} \sum_{p=\lceil n/2\rceil}^{(m+n)/2} \binom{2p}{n} (2^{2p+1}-1) \frac{\zeta(2p+1)}{\pi^{2p}} [u^{m+n-2p}](u\cot u)^{m+1}, \] which is \eqref{eq:general-tangent-power-integral-formula}. \end{proof}

\begin{example}
We illustrate \cref{prop:general-tangent-power-integral}. For \(m=7\), one obtains
\begin{align*}
\int_0^\infty \frac{\tanh^8 x}{x^2}\,dx
&=
\frac{352}{15}\frac{\zeta(3)}{\pi^2}
-\frac{5456}{15}\frac{\zeta(5)}{\pi^4}
+2032\frac{\zeta(7)}{\pi^6}
-4088\frac{\zeta(9)}{\pi^8},
\\
\int_0^\infty \frac{\tanh^8 x}{x^4}\,dx
&=
-\frac{21824}{105}\frac{\zeta(5)}{\pi^4}
+\frac{22352}{3}\frac{\zeta(7)}{\pi^6}
-\frac{228928}{3}\frac{\zeta(9)}{\pi^8}
+245640\frac{\zeta(11)}{\pi^{10}},
\\
\int_0^\infty \frac{\tanh^8 x}{x^6}\,dx
&=
\frac{44704}{35}\frac{\zeta(7)}{\pi^6}
-\frac{1259104}{15}\frac{\zeta(9)}{\pi^8}
+1375584\frac{\zeta(11)}{\pi^{10}}
-6487272\frac{\zeta(13)}{\pi^{12}},
\\
\int_0^\infty \frac{\tanh^8 x}{x^8}\,dx
&=
-\frac{102784}{15}\frac{\zeta(9)}{\pi^8}
+720544\frac{\zeta(11)}{\pi^{10}}
-17299392\frac{\zeta(13)}{\pi^{12}}
+112456344\frac{\zeta(15)}{\pi^{14}}.
\end{align*}
In particular, as \(n\) increases, the first zeta values disappear:
for \(n=3\) there is no \(\zeta(3)\)-term, for \(n=5\) there are no
\(\zeta(3)\)- and \(\zeta(5)\)-terms, and for \(n=7\) the
\(\zeta(3)\)-, \(\zeta(5)\)-, and \(\zeta(7)\)-terms are absent.
\end{example}

\begin{example}
For \(m=8\), the same formula gives
\begin{align*}
\int_0^\infty \frac{\tanh^9 x}{x^3}\,dx
&=
-7\frac{\zeta(3)}{\pi^2}
+\frac{50716}{105}\frac{\zeta(5)}{\pi^4}
-7239\frac{\zeta(7)}{\pi^6}
+42924\frac{\zeta(9)}{\pi^8}
-92115\frac{\zeta(11)}{\pi^{10}},
\\
\int_0^\infty \frac{\tanh^9 x}{x^5}\,dx
&=
31\frac{\zeta(5)}{\pi^4}
-\frac{103886}{21}\frac{\zeta(7)}{\pi^6}
+135926\frac{\zeta(9)}{\pi^8}
-1289610\frac{\zeta(11)}{\pi^{10}}
+4054545\frac{\zeta(13)}{\pi^{12}},
\\
\int_0^\infty \frac{\tanh^9 x}{x^7}\,dx
&=
-127\frac{\zeta(7)}{\pi^6}
+\frac{1671992}{45}\frac{\zeta(9)}{\pi^8}
-1633506\frac{\zeta(11)}{\pi^{10}}
+22705452\frac{\zeta(13)}{\pi^{12}}
-98399301\frac{\zeta(15)}{\pi^{14}},
\\
\int_0^\infty \frac{\tanh^9 x}{x^9}\,dx
&=
511\frac{\zeta(9)}{\pi^8}
-\frac{1674446}{7}\frac{\zeta(11)}{\pi^{10}}
+15407271\frac{\zeta(13)}{\pi^{12}}
-295197903\frac{\zeta(15)}{\pi^{14}}
+1686883770\frac{\zeta(17)}{\pi^{16}}.
\end{align*}
Again the lower zeta values are absent exactly as predicted by the lower
summation bound \(p=\lceil n/2\rceil\), or equivalently by the vanishing
of \(\displaystyle\binom{2p}{n}\) for \(2p<n\).
\end{example}

\begin{remark}
With the change of notation \(a=m+1\) and \(b=n+1\), the integral in
Proposition~\ref{prop:general-tangent-power-integral} belongs to the family
\[
K(a,b):=\int_0^\infty \frac{\tanh^a x}{x^b}\,dx,
\qquad a\geq b\geq 2,\qquad a\equiv b \pmod 2.
\]
This family is known to admit closed forms in terms of odd zeta values, often
written in umbral notation involving Mittag--Leffler polynomials.%
\footnote{For background discussions of this integral family, see
\url{https://mathoverflow.net/questions/271526/is-there-a-closed-form-for-int-0-infty-frac-tanh3xx2dx}
and \url{https://en.wikipedia.org/wiki/Mittag-Leffler_polynomials\#Closed_forms_of_integral_families}.}
The point of formula~\eqref{eq:general-tangent-power-integral-formula} is that it
gives an explicit coefficient form, with the coefficients extracted directly from
\((u\cot u)^{m+1}\), rather than using umbral notation.
\end{remark}

\begin{remark}[The last zeta coefficient]
The last summand in
\eqref{eq:general-tangent-power-integral-formula} corresponds to
\[
p=\frac{m+n}{2}.
\]
For this value one has
\[
m+n-2p=0,
\]
and therefore
\[
[u^0](u\cot u)^{m+1}=1.
\]
Hence the highest zeta term is always
\[
(-1)^{(m-n)/2}
\binom{m+n}{n}
(2^{m+n+1}-1)
\frac{\zeta(m+n+1)}{\pi^{m+n}}.
\]
In particular, the coefficient of the last zeta value is an integer.
\end{remark}

\begin{remark}[The first zeta coefficient for even \(n\)] Assume that \(n\) is even. Since \(m+n\) is even, \(m\) is even as well. The first summand in \eqref{eq:general-tangent-power-integral-formula} corresponds to \(p=n/2\). For this value, \[ \binom{2p}{n}=\binom{n}{n}=1 \] and \[ [u^{m+n-2p}](u\cot u)^{m+1} = [u^m](u\cot u)^{m+1}. \] By the Lagrange inversion  formula \textnormal{\cite{SuryaWarnke2023},\cite[Theorem~5.4.2]{Stanley2023}; see also \cite{Gessel2016}}, applied to the compositional inverse pair \[ u=\arctan t, \qquad t=\tan u, \] we obtain \[ [u^m]\left(\frac{u}{\tan u}\right)^{m+1} = (m+1)[t^{m+1}]\arctan t. \] Since \(m\) is even and \[ [t^{m+1}]\arctan t = \frac{(-1)^{m/2}}{m+1}, \] it follows that \[ [u^m](u\cot u)^{m+1} = (-1)^{m/2}. \] Consequently, the first term in \eqref{eq:general-tangent-power-integral-formula} is \[ (-1)^{(m-n)/2}(-1)^{m/2} (2^{n+1}-1)\frac{\zeta(n+1)}{\pi^n}. \] Because \(m\) is even, \[ (-1)^{(m-n)/2+m/2} = (-1)^{m-n/2} = (-1)^{n/2}, \] and hence the first zeta term is \[ (-1)^{n/2} (2^{n+1}-1)\frac{\zeta(n+1)}{\pi^n}. \] In particular, for even \(n\), the coefficient of the first zeta value is independent of \(m\). \end{remark}

\begin{remark}[Alternation of the relevant coefficients] The alternating signs in the zeta expansion are governed by the coefficients of \((u\cot u)^{m+1}\) occurring in \eqref{eq:general-tangent-power-integral-formula}. These are precisely \[ [u^{2j}](u\cot u)^{m+1}, \qquad 0\leq 2j\leq m. \] Indeed, in \eqref{eq:general-tangent-power-integral-formula}, \[ 2j=m+n-2p, \] and the range of \(p\) implies \(0\leq 2j\leq m\). By the Lagrange inversion formula \textnormal{\cite{SuryaWarnke2023}; see also\cite[Theorem~5.4.2]{Stanley2023} and \cite{Gessel2016,Comtet1974}}, applied to the compositional inverse pair \[ u=\arctan t, \qquad t=\tan u, \] one obtains, throughout this range, \[ [u^{2j}](u\cot u)^{m+1} = \frac{m+1}{m+1-2j} [t^{2j}] \left(\frac{\arctan t}{t}\right)^{m+1-2j}. \] The exponent \[ q:=m+1-2j \] is a positive integer. To determine the sign of the coefficient, use the integral representation \[ \frac{\arctan t}{t} = \int_0^1\frac{dx}{1+t^2x^2}. \] It follows that \[ \left(\frac{\arctan t}{t}\right)^q = \int_{[0,1]^q} \prod_{\ell=1}^q \frac{dx_1\cdots dx_q}{1+t^2x_\ell^2}. \] Expanding each factor as a geometric series gives \[ [t^{2j}] \left(\frac{\arctan t}{t}\right)^q = (-1)^j \int_{[0,1]^q} h_j(x_1^2,\ldots,x_q^2)\, dx_1\cdots dx_q, \] where \(h_j\) denotes the complete homogeneous symmetric polynomial of degree \(j\). Since \[ h_j(x_1^2,\ldots,x_q^2)>0 \] on a set of positive measure in \([0,1]^q\), the integral is strictly positive. Therefore \[ \operatorname{sgn} \left( [t^{2j}] \left(\frac{\arctan t}{t}\right)^q \right) = (-1)^j, \] and hence \[ \operatorname{sgn} \left( [u^{2j}](u\cot u)^{m+1} \right) = (-1)^j. \] Since \[ j=\frac{m+n}{2}-p, \] increasing \(p\) by one changes the sign of the corresponding coefficient. All remaining factors in \eqref{eq:general-tangent-power-integral-formula}, apart from the fixed global sign \((-1)^{(m-n)/2}\), are positive. Consequently, the nonzero zeta coefficients alternate in sign. No assertion is made about the coefficients of \((u\cot u)^{m+1}\) outside the range \(0\leq 2j\leq m\). \end{remark}

\section{Explicit formulae for the Coefficients in Blagouchine's
Generalizations}
\label{sec:blagouchine-coefficients}

For \(a>0\), \(b>0\), and \(r\in\mathbb{N}\), consider the family
of logarithmic integrals
\[
\mathcal{I}_{r}(a,b)
:=
\int_{0}^{\infty}
\frac{\ln(x^{2}+a^{2})}{\cosh^{r}(bx)}\,dx.
\]

Blagouchine evaluated the first several members of this family and
observed that the general form of the answer depends on the parity of
\(r\) \cite{Blagouchine2014}. More precisely, he introduced rational
coefficients \(A_r\) and \(B_{r,\ell}\) in the corresponding odd and
even formulae. However, no closed formulae in terms of \(r\) and
\(\ell\) were given for the complete families \(A_r\) and
\(B_{r,\ell}\). Instead, the coefficients were described as rational
numbers whose values could be computed for each fixed \(r\), and a
table of values up to \(r=20\) was provided.

The coefficient-extraction notation
\[
[w^k]f(w)
\]
will denote the coefficient of \(w^k\) in the Taylor expansion of
\(f(w)\) at \(w=0\). We also write
\[
\Psi_q(z)
:=
\frac{d^q}{dz^q}\Psi(z),
\qquad
\Psi(z)=\frac{\Gamma'(z)}{\Gamma(z)},
\]
so that \(\Psi_0=\Psi\).

The coefficient-extraction formulae below make all the coefficients
in Blagouchine's two parity-dependent expressions explicit.

\subsection{Even powers}

Let \(n\geq 1\). Then
\[
\boxed{
\begin{aligned}
\int_{0}^{\infty}
\frac{\ln(x^{2}+a^{2})}{\cosh^{2n}(bx)}\,dx
={}&
\frac{1}{b}
\frac{
2^{2n-1}\bigl((n-1)!\bigr)^2
}{
(2n-1)!
}
\ln\frac{\pi}{b}
\\
&+
\frac{2}{b}
\sum_{m=0}^{n-1}
\frac{
(-1)^m
[w^{2m}]
\left(\dfrac{w}{\sinh w}\right)^{2n}
}{
(2n-1-2m)!\,
\pi^{\,2n-2-2m}
}
\Psi_{2n-2-2m}\left(\frac12+\frac{ab}{\pi}\right).
\end{aligned}
}
\tag{E}
\]

Equivalently, the coefficients appearing in Blagouchine's notation
are
\[
\boxed{
A_{2n}
=
2(-1)^{n-1}
[w^{2n-2}]
\left(\frac{w}{\sinh w}\right)^{2n}=
\frac{
2^{2n-1}\bigl((n-1)!\bigr)^2
}{
(2n-1)!
}}
\]
and
\[
\boxed{
B_{2n,\ell}
=
\frac{
2(-1)^{n-1-\ell}
}{
(2\ell+1)!
}
[w^{2n-2-2\ell}]
\left(\frac{w}{\sinh w}\right)^{2n},
\qquad
1\leq \ell\leq n-1.
}
\]

\subsection{Odd powers}

Let \(n\geq 0\), and define \(\displaystyle
\xi:=\frac14+\frac{ab}{2\pi},
\qquad
\eta:=\frac34+\frac{ab}{2\pi}.
\)
Then
\[
\boxed{
\begin{aligned}
\int_{0}^{\infty}
\frac{\ln(x^{2}+a^{2})}{\cosh^{2n+1}(bx)}\,dx
={}&
\frac{\pi(2n)!}{4^n b\,(n!)^2}
\ln\left[
\frac{2\pi}{b}
\left(
\frac{\Gamma(\eta)}{\Gamma(\xi)}
\right)^2
\right]
\\
&+
\frac{1}{b}
\sum_{m=0}^{n-1}
\frac{
(-1)^m
[w^{2m}]
\left(\dfrac{w}{\sinh w}\right)^{2n+1}
}{
(2n-2m)!\,
(2\pi)^{\,2n-2m-1}
}
\\
&\hspace{32mm}\times
\left\{
\Psi_{2n-2m-1}(\eta)
-
\Psi_{2n-2m-1}(\xi)
\right\}.
\end{aligned}
}
\tag{O}
\]

The finite sum in \((O)\) is understood to be empty when \(n=0\).

Equivalently, Blagouchine's coefficients in the odd case are
\[
\boxed{
A_{2n+1}
=
2(-1)^n
[w^{2n}]
\left(\frac{w}{\sinh w}\right)^{2n+1}=
2\,\frac{\binom{2n}{n}}{4^n}
}
\]
and
\[
\boxed{
B_{2n+1,\ell}
=
\frac{
(-1)^{n-\ell}
}{
2^{\,2\ell-1}(2\ell)!
}
[w^{2n-2\ell}]
\left(\frac{w}{\sinh w}\right)^{2n+1},
\qquad
1\leq \ell\leq n.
}
\]

Thus, formulae \((E)\) and \((O)\) replace the previously unspecified
rational coefficients by explicit coefficient extractions from the
single elementary generating function
\(\displaystyle
\left(\frac{w}{\sinh w}\right)^r.
\)

\section{Malmsten-type integrals for $\boldsymbol{\beta(2n)}$ and
$\boldsymbol{\zeta(2n+1)}$}
\label{sec:Malmsten-integrals}

Malmsten-type integrals involving the logarithmic kernel
$\ln\ln x$ have long been used to represent special values of
Dirichlet series. A classical example is the evaluation of
$\beta(2)$ obtained by Adamchik \cite{Adamchik1997},
\begin{equation}\label{eq:Adamchik-beta2}
\beta(2)
=
\frac{\pi}{2}
\int_{1}^{\infty}
\frac{x^{4}-6x^{2}+1}{(x^{2}+1)^{3}}
\ln(\ln x)\,dx.
\end{equation}
Similarly, Blagouchine \cite{Blagouchine2014} established the
representation
\begin{equation}\label{eq:Blagouchine-zeta3}
\zeta(3)
=
\frac{8\pi^{2}}{7}
\int_{1}^{\infty}
\frac{x(x^{4}-4x^{2}+1)}{(x^{2}+1)^{4}}
\ln(\ln x)\,dx.
\end{equation}

More recently, Kyrion \cite{Kyrion2025} generalized these isolated
evaluations to the complete families of even beta values and odd
zeta values. In particular, for $N\in\mathbb N$,
\begin{align}
\beta(2N)
&=
(-1)^{N}
\frac{\pi^{2N-1}}
     {2^{2N-1}(2N-1)!}
\int_{1}^{\infty}
\sum_{k=0}^{N-1}
\Bigg(
\sum_{j=0}^{k}
\binom{2k+1}{k-j}
(-1)^{j+1}(2j+1)^{2N-1}
\nonumber\\
&\hspace{5cm}\times
\frac{
\bigl((2k+1)x^{4}-2(2k+3)x^{2}+2k+1\bigr)x^{2k}
}
{(x^{2}+1)^{2k+3}}
\Bigg)
\ln(\ln x)\,dx,
\label{eq:Kyrion-beta}
\end{align}
and
\begin{align}
\zeta(2N+1)
&=
(-1)^{N}
\frac{2\pi^{2N}}
     {(2^{2N+1}-1)(2N)!}
\int_{1}^{\infty}
\sum_{k=1}^{N}
\Bigg(
\sum_{j=1}^{k}
\binom{2k}{k-j}
(-1)^{j}(2j)^{2N}
\nonumber\\
&\hspace{5cm}\times
\frac{
\bigl(2kx^{4}-2(2k+2)x^{2}+2k\bigr)x^{2k-1}
}
{(x^{2}+1)^{2k+2}}
\Bigg)
\ln(\ln x)\,dx.
\label{eq:Kyrion-zeta}
\end{align}

Thus, \cref{eq:Adamchik-beta2,eq:Blagouchine-zeta3} appear as the
first instances of general families. However, the price paid for the
generality of \cref{eq:Kyrion-beta,eq:Kyrion-zeta} is the occurrence
of nested finite sums, which considerably obscures the structure of
the corresponding kernels.

A much more compact representation was subsequently obtained in
\cite{TallaWaffo2025arxiv2511.02843}. There, the same two families
were written in terms of negative-order polylogarithms as
\begin{equation}\label{eq:compact-polylogarithmic-representations}
\begin{dcases}
\displaystyle
\int_{0}^{1}
\frac{\Im\!\left(\operatorname{Li}_{-2n}(ix)\right)}{x}
\ln\!\ln\frac{1}{x}\,dx
=
(-1)^{n+1}
\frac{2^{2n-1}(2n-1)!}{\pi^{2n-1}}
\beta(2n),
\\[3ex]
\displaystyle
\int_{0}^{1}
\frac{\operatorname{Li}_{-2n-1}(-x^{2})}{x}
\ln\!\ln\frac{1}{x}\,dx
=
(-1)^n
\left(1-2^{-(2n+1)}\right)
\frac{(2n)!}{2\pi^{2n}}
\zeta(2n+1).
\end{dcases}
\end{equation}

Compared with \cref{eq:Kyrion-beta,eq:Kyrion-zeta},
\cref{eq:compact-polylogarithmic-representations} eliminates the
nested combinatorial sums altogether and encodes the corresponding
Malmsten kernels in the two compact functions
$\Im\operatorname{Li}_{-2n}(ix)$ and
$\operatorname{Li}_{-2n-1}(-x^{2})$.

The proof given in \cite{TallaWaffo2025arxiv2511.02843}, however,
relied on interchanges of infinite summation and integration whose
validity was not established in that argument. The aim of the present
section is therefore to give a new and rigorous proof of
\cref{eq:compact-polylogarithmic-representations}, avoiding such
unjustified sum--integral interchanges. Our approach instead relies
on finite Eulerian expansions, shifted hyperbolic integrals,
Chebyshev polynomials of the second kind, and a coefficient-collapse
phenomenon.

Throughout this section, we use the notation 
\(\displaystyle\left\langle {m\atop k}\right\rangle\) for Eulerian numbers of type \(A\) and 
\(\displaystyle\left\langle {m\atop k}\right\rangle^{B}\) for Eulerian numbers of type \(B\).

\begin{lemma}[Eulerian--Chebyshev sums]
\label[lemma]{lem:Eulerian-Chebyshev-sums}
Let $n\geq 1$. Then the following identities hold as identities of meromorphic functions of $t \in \mathbb{C}$.
\begin{equation}\label{eq:Eulerian-Chebyshev-sums}
\begin{dcases}
\displaystyle
\sum_{k=0}^{n-1}
(-1)^k
\left\langle {2n-1\atop k}\right\rangle^{B}
U_{2n-2k-2}(\cosh t)
=
-2^{2n-2}
\frac{\cosh^{2n}t}{\sinh t}
\frac{d^{\,2n-1}}{dt^{\,2n-1}}
\operatorname{sech}t,
\\[2ex]
\displaystyle
\sum_{k=0}^{n-1}
(-1)^k
\left\langle {2n\atop k}\right\rangle
U_{2n-2k-2}(\cosh t)
=
-\frac{1}{2}
\frac{\cosh^{2n+1}t}{\sinh t}
\frac{d^{\,2n}}{dt^{\,2n}}
\tanh t.
\end{dcases}
\end{equation}
Here
$\left\langle {m\atop k}\right\rangle$ and
$\left\langle {m\atop k}\right\rangle^{B}$
denote the Eulerian numbers of type $A$ and type $B$, respectively,
and $U_m$ denotes the Chebyshev polynomial of the second kind.
\end{lemma}

\begin{proof}
From the hyperbolic representation
\[
U_m(\cosh t)
=
\frac{\sinh((m+1)t)}{\sinh t},
\]
we obtain
\begin{equation}\label{eq:Euler-Cheb-hyperbolic}
\begin{dcases}
\displaystyle
\sum_{k=0}^{n-1}
(-1)^k
\left\langle {2n-1\atop k}\right\rangle^{B}
U_{2n-2k-2}(\cosh t)
=
\frac{1}{\sinh t}
\sum_{k=0}^{n-1}
(-1)^k
\left\langle {2n-1\atop k}\right\rangle^{B}
\sinh((2n-2k-1)t),
\\[2ex]
\displaystyle
\sum_{k=0}^{n-1}
(-1)^k
\left\langle {2n\atop k}\right\rangle
U_{2n-2k-2}(\cosh t)
=
\frac{1}{\sinh t}
\sum_{k=0}^{n-1}
(-1)^k
\left\langle {2n\atop k}\right\rangle
\sinh((2n-2k-1)t).
\end{dcases}
\end{equation}

Using
\[
\sinh x=\frac{e^x-e^{-x}}{2},
\]
the sums on the right-hand side become
\begin{equation}\label{eq:Euler-Cheb-exponential}
\begin{dcases}
\displaystyle
\sum_{k=0}^{n-1}
(-1)^k
\left\langle {2n-1\atop k}\right\rangle^{B}
\sinh((2n-2k-1)t)
\\
\displaystyle\qquad
=
\frac{e^{(2n-1)t}}{2}
\sum_{k=0}^{n-1}
(-1)^k
\left\langle {2n-1\atop k}\right\rangle^{B}
\left(
e^{-2kt}
-
e^{-2(2n-1-k)t}
\right),
\\[3ex]
\displaystyle
\sum_{k=0}^{n-1}
(-1)^k
\left\langle {2n\atop k}\right\rangle
\sinh((2n-2k-1)t)
\\
\displaystyle\qquad
=
\frac{e^{(2n-1)t}}{2}
\sum_{k=0}^{n-1}
(-1)^k
\left\langle {2n\atop k}\right\rangle
\left(
e^{-2kt}
-
e^{-2(2n-1-k)t}
\right).
\end{dcases}
\end{equation}

We now use the symmetry relations \cite{FulmanLeeKimPetersen2021}
\[
\left\langle{m\atop r}\right\rangle=\left\langle{m\atop m-1-r}\right\rangle,
\qquad
\left\langle{m\atop r}\right\rangle^{B}=\left\langle{m\atop m-r}\right\rangle^{B}.
\]
In the second part of each finite sum, set \(j=2n-1-k.\) Therefore,
\begin{equation}
\begin{dcases}
\displaystyle
-\sum_{k=0}^{n-1}(-1)^k
\left\langle {2n-1\atop k}\right\rangle^{B}
e^{-2(2n-1-k)t}
=
\sum_{j=n}^{2n-1}(-1)^j
\left\langle {2n-1\atop j}\right\rangle^{B}
e^{-2jt},
\\[2ex]
\displaystyle
-\sum_{k=0}^{n-1}(-1)^k
\left\langle {2n\atop k}\right\rangle
e^{-2(2n-1-k)t}
=
\sum_{j=n}^{2n-1}(-1)^j
\left\langle {2n\atop j}\right\rangle
e^{-2jt}.
\end{dcases}
\end{equation}

Consequently,
\begin{equation}\label{eq:Euler-Cheb-completed}
\begin{dcases}
\displaystyle
\sum_{k=0}^{n-1}
(-1)^k
\left\langle {2n-1\atop k}\right\rangle^{B}
\left(
e^{-2kt}
-
e^{-2(2n-1-k)t}
\right)
=
\sum_{k=0}^{2n-1}
(-1)^k
\left\langle {2n-1\atop k}\right\rangle^{B}
e^{-2kt},
\\[3ex]
\displaystyle
\sum_{k=0}^{n-1}
(-1)^k
\left\langle {2n\atop k}\right\rangle
\left(
e^{-2kt}
-
e^{-2(2n-1-k)t}
\right)
=
\sum_{k=0}^{2n-1}
(-1)^k
\left\langle {2n\atop k}\right\rangle
e^{-2kt}.
\end{dcases}
\end{equation}

Introduce the Eulerian polynomials
\[
B_m(x)
=
\sum_{k=0}^{m}
\left\langle {m\atop k}\right\rangle^{B}x^k,
\qquad
A_m(x)
=
\sum_{k=0}^{m-1}
\left\langle {m\atop k}\right\rangle x^k.
\]
Equation \eqref{eq:Euler-Cheb-completed} then gives
\begin{equation}\label{eq:Euler-Cheb-polynomials}
\begin{dcases}
\displaystyle
\sum_{k=0}^{n-1}
(-1)^k
\left\langle {2n-1\atop k}\right\rangle^{B}
U_{2n-2k-2}(\cosh t)
=
\frac{e^{(2n-1)t}}{2\sinh t}
B_{2n-1}(-e^{-2t}),
\\[3ex]
\displaystyle
\sum_{k=0}^{n-1}
(-1)^k
\left\langle {2n\atop k}\right\rangle
U_{2n-2k-2}(\cosh t)
=
\frac{e^{(2n-1)t}}{2\sinh t}
A_{2n}(-e^{-2t}).
\end{dcases}
\end{equation}

We now employ the classical generating identities
\[
\frac{B_m(x)}{(1-x)^{m+1}}
=
\sum_{r=0}^{\infty}(2r+1)^m x^r,
\qquad
\frac{A_m(x)}{(1-x)^{m+1}}
=
\sum_{r=0}^{\infty}(r+1)^m x^r.
\]
Taking $x=-e^{-2t}$ with $\Re t > 0$ yields
\begin{equation}\label{eq:Euler-Cheb-generating}
\begin{dcases}
\displaystyle
B_{2n-1}(-e^{-2t})
=
(1+e^{-2t})^{2n}
\sum_{r=0}^{\infty}
(-1)^r(2r+1)^{2n-1}e^{-2rt},
\\[3ex]
\displaystyle
A_{2n}(-e^{-2t})
=
(1+e^{-2t})^{2n+1}
\sum_{r=0}^{\infty}
(-1)^r(r+1)^{2n}e^{-2rt}.
\end{dcases}
\end{equation}

Since \(\displaystyle1+e^{-2t}=2e^{-t}\cosh t,\) substitution into \eqref{eq:Euler-Cheb-polynomials} gives
\begin{equation}\label{eq:Euler-Cheb-series}
\begin{dcases}
\displaystyle
\sum_{k=0}^{n-1}
(-1)^k
\left\langle {2n-1\atop k}\right\rangle^{B}
U_{2n-2k-2}(\cosh t)
=
2^{2n-1}
\frac{\cosh^{2n}t}{\sinh t}
\sum_{r=0}^{\infty}
(-1)^r(2r+1)^{2n-1}e^{-(2r+1)t},
\\[3ex]
\displaystyle
\sum_{k=0}^{n-1}
(-1)^k
\left\langle {2n\atop k}\right\rangle
U_{2n-2k-2}(\cosh t)
=
2^{2n}
\frac{\cosh^{2n+1}t}{\sinh t}
\sum_{r=1}^{\infty}
(-1)^{r-1}r^{2n}e^{-2rt}.
\end{dcases}
\end{equation}

Finally,
\[
\operatorname{sech}t
=
\frac{2e^{-t}}{1+e^{-2t}}
=
2\sum_{r=0}^{\infty}
(-1)^r e^{-(2r+1)t},
\]
and therefore
\[
\frac{d^{\,2n-1}}{dt^{\,2n-1}}
\operatorname{sech}t
=
-2\sum_{r=0}^{\infty}
(-1)^r(2r+1)^{2n-1}e^{-(2r+1)t}.
\]
Moreover,
\[
\sum_{r=1}^{\infty}
(-1)^{r-1}e^{-2rt}
=
\frac{1}{e^{2t}+1}
=
\frac{1-\tanh t}{2},
\]
so that
\[
2^{2n}
\sum_{r=1}^{\infty}
(-1)^{r-1}r^{2n}e^{-2rt}
=
-\frac{1}{2}
\frac{d^{\,2n}}{dt^{\,2n}}
\tanh t.
\]
Substituting these two identities into
\eqref{eq:Euler-Cheb-series} proves
\eqref{eq:Eulerian-Chebyshev-sums}. Since both sides are meromorphic functions of t, the identities extend to $\mathbb{C}$ by analytic continuation.
\end{proof}

\begin{lemma}[Eulerian--Chebyshev coefficient collapse]
\label[lemma]{lem:Eulerian-Chebyshev-collapse}
Let $n\geq 1$ and $p\geq 1$. Then
\begin{equation}\label{eq:Eulerian-Chebyshev-collapse}
\begin{dcases}
\displaystyle
[z^{2n-1}]\left\{
(\arcsin z)^{2p-1}
\sum_{k=0}^{n-1}
(-1)^k
\left\langle {2n-1\atop k}\right\rangle^{B}
U_{2n-2k-2}(z)\right\}
=
(-1)^{n-1}2^{2n-2}(2n-1)!\,\delta_{p,n},
\\[3ex]
\displaystyle
[z^{2n}]\left\{
(\arcsin z)^{2p}
\sum_{k=0}^{n-1}
(-1)^k
\left\langle {2n\atop k}\right\rangle
U_{2n-2k-2}(z)\right\}
=
\frac{(-1)^{n-1}}{2}(2n)!\,\delta_{p,n},
\end{dcases}
\end{equation}
where $\delta_{p,n}$ denotes the Kronecker delta.
\end{lemma}

\begin{proof}
We first transform the identities of
\cref{lem:Eulerian-Chebyshev-sums}.
Set \(\displaystyle
t=i\left(\frac{\pi}{2}-y\right).\)
Then
\[
\cosh t=\sin y,
\qquad
\sinh t=i\cos y,
\qquad
\operatorname{sech}t=\csc y,
\qquad
\tanh t=i\cot y,
\]
and
\[
\frac{d}{dt}=i\frac{d}{dy}.
\]
Hence \cref{lem:Eulerian-Chebyshev-sums} gives
\begin{equation}\label{eq:Eulerian-Chebyshev-trigonometric}
\begin{dcases}
\displaystyle
\sum_{k=0}^{n-1}
(-1)^k
\left\langle {2n-1\atop k}\right\rangle^{B}
U_{2n-2k-2}(\sin y)
=
(-1)^n2^{2n-2}
\frac{\sin^{2n}y}{\cos y}
\frac{d^{\,2n-1}}{dy^{\,2n-1}}\csc y,
\\[3ex]
\displaystyle
\sum_{k=0}^{n-1}
(-1)^k
\left\langle {2n\atop k}\right\rangle
U_{2n-2k-2}(\sin y)
=
\frac{(-1)^{n-1}}{2}
\frac{\sin^{2n+1}y}{\cos y}
\frac{d^{\,2n}}{dy^{\,2n}}\cot y.
\end{dcases}
\end{equation}

Let $C$ be a sufficiently small positively oriented circle
centered at the origin, contained in the disk of analyticity of
the principal branch of $\arcsin$. By Cauchy's coefficient formula,
\[
[z^m]F(z)
=
\frac{1}{2\pi i}
\int_C
\frac{F(z)}{z^{m+1}}\,dz.
\]

For the first coefficient in
\cref{eq:Eulerian-Chebyshev-collapse}, we therefore have
\begin{align*}
&[z^{2n-1}]\left\{
(\arcsin z)^{2p-1}
\sum_{k=0}^{n-1}
(-1)^k
\left\langle {2n-1\atop k}\right\rangle^{B}
U_{2n-2k-2}(z)\right\}
\\
&\qquad=
\frac{1}{2\pi i}
\int_C
\frac{(\arcsin z)^{2p-1}}{z^{2n}}
\sum_{k=0}^{n-1}
(-1)^k
\left\langle {2n-1\atop k}\right\rangle^{B}
U_{2n-2k-2}(z)\,dz.
\end{align*}

Make the substitution
\[
z=\sin y,
\qquad
dz=\cos y\,dy,
\]
and let $\Gamma=\arcsin(C)$ be the corresponding positively
oriented contour around the origin. Using
\cref{eq:Eulerian-Chebyshev-trigonometric}, we obtain
\begin{align*}
&[z^{2n-1}]\left\{
(\arcsin z)^{2p-1}
\sum_{k=0}^{n-1}
(-1)^k
\left\langle {2n-1\atop k}\right\rangle^{B}
U_{2n-2k-2}(z)\right\}
\\
&\qquad=
\frac{(-1)^n2^{2n-2}}{2\pi i}
\int_{\Gamma}
y^{2p-1}
\frac{d^{\,2n-1}}{dy^{\,2n-1}}\csc y\,dy.
\end{align*}

Integrating by parts $2n-1$ times along the closed contour
$\Gamma$ gives
\begin{align}
&[z^{2n-1}]\left\{
(\arcsin z)^{2p-1}
\sum_{k=0}^{n-1}
(-1)^k
\left\langle {2n-1\atop k}\right\rangle^{B}
U_{2n-2k-2}(z)\right\}
\nonumber\\
&\qquad=
\frac{(-1)^{n-1}2^{2n-2}}{2\pi i}
\int_{\Gamma}
\frac{d^{\,2n-1}}{dy^{\,2n-1}}
\left(y^{2p-1}\right)
\csc y\,dy.
\label{eq:collapse-B-Cauchy}
\end{align}

We now distinguish three cases:
\[
\frac{d^{\,2n-1}}{dy^{\,2n-1}}
\left(y^{2p-1}\right)
=
\begin{cases}
0,
& p<n,
\\[1ex]
(2n-1)!,
& p=n,
\\[1ex]
\displaystyle
\frac{(2p-1)!}{(2p-2n)!}\,
y^{2p-2n},
& p>n.
\end{cases}
\]

If $p<n$, the integral in
\cref{eq:collapse-B-Cauchy} vanishes immediately.

If $p=n$, then
\[
\frac{1}{2\pi i}
\int_{\Gamma}\csc y\,dy
=
\frac{1}{2\pi i}
\int_{\Gamma}
\frac{1}{y}\frac{y}{\sin y}\,dy
=1,
\]
by Cauchy's integral formula, since $y/\sin y$ is holomorphic
near $0$ and \(\displaystyle\left.\frac{y}{\sin y}\right|_{y=0}=1.\)
Thus
\[
[z^{2n-1}]\left\{
(\arcsin z)^{2n-1}
\sum_{k=0}^{n-1}
(-1)^k
\left\langle {2n-1\atop k}\right\rangle^{B}
U_{2n-2k-2}(z)\right\}
=
(-1)^{n-1}2^{2n-2}(2n-1)!.
\]

If $p>n$, then
\[
y^{2p-2n}\csc y
=
y^{2p-2n-1}\frac{y}{\sin y}
\]
is holomorphic inside $\Gamma$, and therefore its contour integral
vanishes by Cauchy's theorem. Hence
\[
[z^{2n-1}]\left\{
(\arcsin z)^{2p-1}
\sum_{k=0}^{n-1}
(-1)^k
\left\langle {2n-1\atop k}\right\rangle^{B}
U_{2n-2k-2}(z)\right\}
=
(-1)^{n-1}2^{2n-2}(2n-1)!\,\delta_{p,n}.
\]

For the second coefficient, Cauchy's coefficient formula gives
\begin{align*}
&[z^{2n}]\left\{
(\arcsin z)^{2p}
\sum_{k=0}^{n-1}
(-1)^k
\left\langle {2n\atop k}\right\rangle
U_{2n-2k-2}(z)\right\}
\\
&\qquad=
\frac{1}{2\pi i}
\int_C
\frac{(\arcsin z)^{2p}}{z^{2n+1}}
\sum_{k=0}^{n-1}
(-1)^k
\left\langle {2n\atop k}\right\rangle
U_{2n-2k-2}(z)\,dz.
\end{align*}

Again setting
\[
z=\sin y,
\qquad
dz=\cos y\,dy,
\]
and using
\cref{eq:Eulerian-Chebyshev-trigonometric}, we obtain
\begin{align*}
&[z^{2n}]\left\{
(\arcsin z)^{2p}
\sum_{k=0}^{n-1}
(-1)^k
\left\langle {2n\atop k}\right\rangle
U_{2n-2k-2}(z)\right\}
\\
&\qquad=
\frac{(-1)^{n-1}}{2}
\frac{1}{2\pi i}
\int_{\Gamma}
y^{2p}
\frac{d^{\,2n}}{dy^{\,2n}}\cot y\,dy.
\end{align*}

Integrating by parts $2n$ times yields
\begin{align}
&[z^{2n}]\left\{
(\arcsin z)^{2p}
\sum_{k=0}^{n-1}
(-1)^k
\left\langle {2n\atop k}\right\rangle
U_{2n-2k-2}(z)\right\}
\nonumber\\
&\qquad=
\frac{(-1)^{n-1}}{2}
\frac{1}{2\pi i}
\int_{\Gamma}
\frac{d^{\,2n}}{dy^{\,2n}}
\left(y^{2p}\right)
\cot y\,dy.
\label{eq:collapse-A-Cauchy}
\end{align}

Here
\[
\frac{d^{\,2n}}{dy^{\,2n}}
\left(y^{2p}\right)
=
\begin{cases}
0,
& p<n,
\\[1ex]
(2n)!,
& p=n,
\\[1ex]
\displaystyle
\frac{(2p)!}{(2p-2n)!}\,
y^{2p-2n},
& p>n.
\end{cases}
\]

If $p<n$, the integral in
\cref{eq:collapse-A-Cauchy} vanishes.

If $p=n$, then
\[
\frac{1}{2\pi i}
\int_{\Gamma}\cot y\,dy
=
\frac{1}{2\pi i}
\int_{\Gamma}
\frac{1}{y}
\frac{y\cos y}{\sin y}\,dy
=1,
\]
since $y\cos y/\sin y$ is holomorphic near $0$ and \(\displaystyle\left.\frac{y\cos y}{\sin y}\right|_{y=0}=1.\)
Therefore
\[
[z^{2n}]\left\{
(\arcsin z)^{2n}
\sum_{k=0}^{n-1}
(-1)^k
\left\langle {2n\atop k}\right\rangle
U_{2n-2k-2}(z)\right\}
=
\frac{(-1)^{n-1}}{2}(2n)!.
\]

Finally, if $p>n$, then
\[
y^{2p-2n}\cot y
=
y^{2p-2n-1}
\frac{y\cos y}{\sin y}
\]
is holomorphic inside $\Gamma$, so its contour integral vanishes
by Cauchy's theorem. Consequently,
\[
[z^{2n}]\left\{
(\arcsin z)^{2p}
\sum_{k=0}^{n-1}
(-1)^k
\left\langle {2n\atop k}\right\rangle
U_{2n-2k-2}(z)\right\}
=
\frac{(-1)^{n-1}}{2}(2n)!\,\delta_{p,n}.
\]

This proves \cref{eq:Eulerian-Chebyshev-collapse}.
\end{proof}

\begin{proposition}[Polylogarithmic integral evaluations]
\label[proposition]{thm:polylogarithmic-integral-evaluations}
For every $n\in\mathbb{N}$,
\begin{equation}\label{eq:polylogarithmic-integral-evaluations}
\begin{dcases}
\displaystyle
\int_{0}^{1}
\frac{\Im\!\left(\operatorname{Li}_{-2n}(ix)\right)}{x}
\ln\!\ln\frac{1}{x}\,dx
=
(-1)^{n-1}2^{2n-1}(2n-1)!
\frac{\beta(2n)}{\pi^{2n-1}},
\\[3ex]
\displaystyle
\int_{0}^{1}
\frac{\operatorname{Li}_{-2n-1}(-x^{2})}{x}
\ln\!\ln\frac{1}{x}\,dx
=
(-1)^n
\frac{(2^{2n+1}-1)(2n)!}{2^{2n+2}}
\frac{\zeta(2n+1)}{\pi^{2n}}.
\end{dcases}
\end{equation}
\end{proposition}

\begin{proof}
We start from the identities established in
\cite{TallaWaffo2026arxiv2602.16761}:
\begin{equation}\label{eq:polylogarithmic-Eulerian-reduction}
\begin{dcases}
\displaystyle
\int_{0}^{1}
\frac{\Im\!\left(\operatorname{Li}_{-2n}(ix)\right)}{x}
\ln\!\ln\frac{1}{x}\,dx
=
-\int_{0}^{1}
\frac{1}{\ln x}\,
\frac{1}{(1+x^{2})^{2n}}
\sum_{k=0}^{n-1}
\left\langle {2n-1\atop k}\right\rangle^{B}
\left\{
(-x^{2})^{k}
+
(-x^{2})^{2n-1-k}
\right\}\,dx,
\\[4ex]
\displaystyle
\int_{0}^{1}
\frac{\operatorname{Li}_{-2n-1}(-x^{2})}{x}
\ln\!\ln\frac{1}{x}\,dx
=
-\frac{1}{2}
\int_{0}^{1}
\frac{1}{x\ln x}\,
\frac{1}{(1+x^{2})^{2n+1}}
\sum_{k=0}^{n-1}
\left\langle {2n\atop k}\right\rangle
\left\{
(-x^{2})^{2n-k}
+
(-x^{2})^{k+1}
\right\}\,dx.
\end{dcases}
\end{equation}

Set \(\displaystyle x=e^{-t}.\) Then
\[
dx=-e^{-t}\,dt,
\qquad
\ln x=-t,
\qquad
1+e^{-2t}=2e^{-t}\cosh t.
\]
Moreover,
\begin{equation}\label{eq:exponential-to-hyperbolic}
\begin{dcases}
\displaystyle
(-e^{-2t})^{k}
+
(-e^{-2t})^{2n-1-k}
=
2(-1)^k e^{-(2n-1)t}
\sinh((2n-2k-1)t),
\\[2ex]
\displaystyle
(-e^{-2t})^{2n-k}
+
(-e^{-2t})^{k+1}
=
-2(-1)^k e^{-(2n+1)t}
\sinh((2n-2k-1)t).
\end{dcases}
\end{equation}
Consequently,
\begin{equation}\label{eq:polylogarithmic-hyperbolic-reduction}
\begin{dcases}
\displaystyle
\int_{0}^{1}
\frac{\Im\!\left(\operatorname{Li}_{-2n}(ix)\right)}{x}
\ln\!\ln\frac{1}{x}\,dx
=
\frac{1}{2^{2n-1}}
\sum_{k=0}^{n-1}
(-1)^k
\left\langle {2n-1\atop k}\right\rangle^{B}
\int_{0}^{\infty}
\frac{\sinh((2n-2k-1)t)}
{t\cosh^{2n}t}\,dt,
\\[4ex]
\displaystyle
\int_{0}^{1}
\frac{\operatorname{Li}_{-2n-1}(-x^{2})}{x}
\ln\!\ln\frac{1}{x}\,dx
=
-\frac{1}{2^{2n+1}}
\sum_{k=0}^{n-1}
(-1)^k
\left\langle {2n\atop k}\right\rangle
\int_{0}^{\infty}
\frac{\sinh((2n-2k-1)t)}
{t\cosh^{2n+1}t}\,dt.
\end{dcases}
\end{equation}

Applying
\cref{prop:shifted_integrals_chebyshev_arcsine}
with \(\displaystyle j=n-k-1,\)
so that
\[
2j+1=2n-2k-1,
\qquad
U_{2j}(z)=U_{2n-2k-2}(z),
\]
and interchanging the finite sums, we obtain
\begin{equation}\label{eq:polylogarithmic-coefficient-reduction}
\begin{dcases}
\displaystyle
\int_{0}^{1}
\frac{\Im\!\left(\operatorname{Li}_{-2n}(ix)\right)}{x}
\ln\!\ln\frac{1}{x}\,dx
=
\sum_{p=1}^{n}
2^{\,2p-2n+1}
\frac{\beta(2p)}{\pi^{2p-1}}
[z^{2n-1}]\left\{
(\arcsin z)^{2p-1}
\sum_{k=0}^{n-1}
(-1)^k
\left\langle {2n-1\atop k}\right\rangle^{B}
U_{2n-2k-2}(z)\right\},
\\[4ex]
\displaystyle
\int_{0}^{1}
\frac{\operatorname{Li}_{-2n-1}(-x^{2})}{x}
\ln\!\ln\frac{1}{x}\,dx
=
-\sum_{p=1}^{n}
\frac{2^{2p+1}-1}{2^{2n+1}}
\frac{\zeta(2p+1)}{\pi^{2p}}
[z^{2n}]\left\{
(\arcsin z)^{2p}
\sum_{k=0}^{n-1}
(-1)^k
\left\langle {2n\atop k}\right\rangle
U_{2n-2k-2}(z)\right\}.
\end{dcases}
\end{equation}

By \cref{lem:Eulerian-Chebyshev-collapse}, the coefficient
expressions in \cref{eq:polylogarithmic-coefficient-reduction}
vanish for every \(p\neq n\). Thus only the term \(p=n\) survives in each sum. This completes the proof.
\end{proof}

\section*{Acknowledgments}

The author acknowledges the use of an AI language model for assistance with the presentation of the manuscript, numerical experimentation, and verification of elementary asymptotic estimates.


\printbibliography

\end{document}
